\section{Spectral characterization}\label{sec:5}
\begingroup
\emergencystretch=3em
\raggedbottom

We will follow Section 5 of [CGSS]. All graphs are finite and simple. For a graph $X$ we write
$V(X)$ and $E(X)$ for its vertices and edges, $A(X)$ for its adjacency matrix and $L(X)$ for
its line graph. The paper's graphs $G$, $G_1$, $G_2$ are our $X$, $X_1$, $X_2$, and its
auxiliary graph $H$ in Theorem 5.8 is our $Y$ (\path{cgss/notation.md}, Section 1). Valencies
are $k$, $k_1$, $k_2$ in Lemmas 5.1 to 5.6 and in Remark 5.7; from Theorem 5.8 on, $k$ is the
block size of a design and $(v,k,b,r,\lambda)$ are the design parameters, as in the paper. A
generic eigenvalue is $\mu$, since $\lambda$ is a design parameter.

The spectrum $\Spec(M)$ of a real symmetric matrix $M$ is the multiset of the roots of its
characteristic polynomial, and $\mult_M(\mu)$ is the number of copies of $\mu$ in it. In a
multiset $\ms{\mu^{m},\dots}$ the exponent $m$ counts copies, and a sum of multisets adds the
counts. The spectrum of a graph is that of its adjacency matrix, two graphs are
\emph{cospectral} when their adjacency matrices have the same characteristic polynomial, and
$\lmax$ is the largest eigenvalue. Where the paper writes that two graphs have the same
spectrum, we will say that they are cospectral. We cite the spectral toolkit \path{cgss/proofs/spectraltools/spectraltools.tex} as [ST],
by the numbers of its results; its structured version \path{cgss/proofs/spectraltools/spectraltools-lamport.tex}
has the same numbering. The toolkit is not part of this document, so we list the facts that we
take from it, with its proposed Lean names, before we use them.

The paper quotes Lemmas 5.1 to 5.4 as well known (errata, Gap 8). The toolkit proves them as
SP1 to SP4, and we will state them in the paper's form and cite the toolkit. Lemmas 5.5 and 5.6
and Remark 5.7 have proofs in the paper, and we will follow those proofs, with the toolkit's
facts at the leaves. For Lemmas 5.1, 5.2, 5.3 and 5.5 the toolkit proposes the same Lean
declarations (its Theorems 5.3, 5.4, 5.6 and 6.5), so each of these results and the
corresponding toolkit result are one declaration. Lemma 5.4, Lemma 5.6 and Remark 5.7 differ in
statement from the toolkit's Theorem 6.4, Theorem 6.7 and Corollary 6.8, in the forms that
Theorem 5.8 needs: Lemma 5.4 does not assume an edge, and asserts $k_1\ne k_2$; Lemma 5.6 (i) holds for
every biregular bipartition of $X_1$ with $|P_1|\le|Q_1|$; and Remark 5.7 assumes only that
$X_2$ has no isolated vertices and is not bipartite. We will derive them from the toolkit's
results, which we cite as inputs, and give them their own Lean names, derived from the
toolkit's by the rules of \path{cgss/notation.md}, Section 10:
\path{SimpleGraph.regular_or_isBiregularWith_of_lineGraph_regular},
\path{SimpleGraph.lineGraph_cospectral_bipartite_forall_isBiregularWith} and
\path{SimpleGraph.lineGraph_cospectral_bipartite_spectrum_of_not_isBipartite}.

\begingroup\raggedright
\medskip\noindent\textbf{Inputs from the spectral toolkit.}\label{input:st}
\repair\ These are the toolkit's proofs of the facts that the paper quotes (errata, Gap 8, and
\path{cgss/prerequisites.md}, SP1 to SP5). Matrices are real, square matrices are symmetric,
and index sets are finite.
\begin{itemize}
 \item \factlabel{[ST, Theorem 2.2]}{st:2.2}\ST{Theorem 2.2}
  (\path{Matrix.IsHermitian.eigenvalueMultiplicity_eq_finrank},
  \path{Matrix.IsHermitian.mem_spectrumMultiset_iff}). If $M$ is symmetric of order $n$, then
  $\Spec(M)$ has $n$ members counted with multiplicity, $\mult_M(\mu)=\dim\ker(M-\mu I)$, and
  $\mu\in\Spec(M)$ exactly when $M\Zx=\mu\Zx$ for some $\Zx\ne0$.
 \item \factlabel{[ST, Lemma 2.4]}{st:2.4}\ST{Lemma 2.4}
  (\path{Matrix.IsHermitian.largestEigenvalue_mem_spectrumMultiset},
  \path{Matrix.IsHermitian.le_largestEigenvalue_of_mem}). If $M$ is symmetric of positive
  order, then $\lmax(M)$ is the largest member of $\Spec(M)$.
 \item \factlabel{[ST, Proposition 2.6]}{st:2.6}\ST{Proposition 2.6}
  (\path{Matrix.IsHermitian.spectrumMultiset_eq_iff_charpoly_eq},
  \path{SimpleGraph.cospectral_iff_spectrumMultiset_eq}, \path{SimpleGraph.Iso.cospectral}).
  For symmetric $M$, $M'$, possibly on different index sets, $\chi_M=\chi_{M'}$ if and only if
  $\Spec(M)=\Spec(M')$, if and only if $\mult_M(\mu)=\mult_{M'}(\mu)$ for every real $\mu$.
  Cospectral graphs have the same number of vertices, and isomorphic graphs are cospectral.
 \item \factlabel{[ST, Theorem 3.1]}{st:3.1}\ST{Theorem 3.1}
  (\path{Matrix.IsHermitian.spectrumMultiset_aeval}). If $M$ is symmetric and $f$ is a real
  polynomial, then $\Spec(f(M))$ is the image of $\Spec(M)$ under $f$, with multiplicities. In
  particular $\Spec(cI+M)$ is $\Spec(M)$ shifted by $c$, and $\mult_{cI+M}(\mu+c)=\mult_M(\mu)$.
 \item \factlabel{[ST, Theorem 3.2]}{st:3.2}\ST{Theorem 3.2}
  (\path{Matrix.IsHermitian.rank_aeval_eq_sum_multiplicity}). If $M$ is symmetric of order $n$,
  then $\rk M=n-\mult_M(0)$, and for a real polynomial $f$, $\rk f(M)$ is the sum of
  $\mult_M(\theta)$ over the distinct eigenvalues $\theta$ of $M$ with $f(\theta)\ne0$.
 \item \factlabel{[ST, Corollary 3.4]}{st:3.4}\ST{Corollary 3.4}
  (\path{SimpleGraph.posSemidef_two_add_adjMatrix_iff}). For a graph $X$, the matrix
  $2I+A(X)$ is positive semidefinite if and only if every member of $\Spec(X)$ is at least $-2$.
 \item \factlabel{[ST, Theorem 3.5]}{st:3.5}\ST{Theorem 3.5}
  (\path{Matrix.rank_gram_eq_finrank_span}). For every matrix $B$,
  $\rk(B^TB)=\rk B=\rk(BB^T)$. The Gram matrix of a finite family of vectors in a real inner
  product space is positive semidefinite, and its rank is the dimension of the span of the family.
 \item \factlabel{[ST, Theorem 5.1]}{st:5.1}\ST{Theorem 5.1}
  (\path{SimpleGraph.adjMatrix_cube_apply_self_eq_twice_triangles}). For a graph $X$ with
  $A=A(X)$ and a vertex $u$, $A^3(u,u)$ is twice the number of triangles of $X$ that
  contain $u$.
 \item \factlabel{[ST, Theorem 5.3]}{st:5.3}\ST{Theorem 5.3} (SP1;
  \path{SimpleGraph.IsRegularOfDegree.largestEigenvalue_eq},
  \path{SimpleGraph.IsRegularOfDegree.multiplicity_eq_card_components}). If $X$ is
  $k$-regular on a nonempty vertex set, then $\lmax(X)=k$ and $\mult_X(k)$ is the number of
  connected components of $X$.
 \item \factlabel{[ST, Theorem 5.4]}{st:5.4}\ST{Theorem 5.4} (SP2;
  \path{SimpleGraph.isRegular_iff_sum_sq_eigenvalues}). A graph $X$ on $n\ge1$ vertices is
  regular if and only if $\sum_{\mu\in\Spec(X)}\mu^2=n\lmax(X)$.
 \item \factlabel{[ST, Corollary 5.5]}{st:5.5}\ST{Corollary 5.5}
  (\path{SimpleGraph.Cospectral.isRegularOfDegree},
  \path{SimpleGraph.Cospectral.connected_of_regular}). If $X$ is $k$-regular on a nonempty
  vertex set and $X'$ is cospectral with $X$, then $X'$ is $k$-regular, and $X'$ is connected
  if $X$ is.
 \item \factlabel{[ST, Theorem 5.6]}{st:5.6}\ST{Theorem 5.6} (SP3;
  \path{SimpleGraph.IsRegularOfDegree.isBipartite_iff_neg_eigenvalue}). If $X$ is connected and
  $k$-regular, then $X$ is bipartite if and only if $-k\in\Spec(X)$, and then $\mult_X(-k)=1$.
 \item \factlabel{[ST, Corollary 5.9]}{st:5.9}\ST{Corollary 5.9}
  (\path{Matrix.spectrumMultiset_smul_one_add_smul_allOnes}). For real $c$, $t$ and order
  $n>0$, $\Spec(cI+tJ)=\ms{(c+tn)^{1},c^{n-1}}$.
 \item \factlabel{[ST, Theorem 6.1]}{st:6.1}\ST{Theorem 6.1}
  (\path{Matrix.eigenvalueMultiplicity_mul_comm_of_ne_zero},
  \path{Matrix.spectrumMultiset_gram_products}). If $M$ is $P\times Q$ and $N$ is $Q\times P$,
  then $\mult_{MN}(\mu)=\mult_{NM}(\mu)$ for every $\mu\ne0$. If $B$ is $P\times Q$ of rank
  $\rho$, then $BB^T$ and $B^TB$ are positive semidefinite, have the same positive eigenvalues
  with multiplicities, and have $\mult_{BB^T}(0)=|P|-\rho$ and $\mult_{B^TB}(0)=|Q|-\rho$.
 \item \factlabel{[ST, Theorem 6.2]}{st:6.2}\ST{Theorem 6.2}
  (\path{SimpleGraph.edgeIncMatrix_mul_transpose}, \path{SimpleGraph.transpose_edgeIncMatrix_mul},
  \path{SimpleGraph.multiplicity_neg_two_lineGraph}). Let $M$ be the $V(X)\times E(X)$
  incidence matrix of a graph $X$ and $D$ its diagonal matrix of degrees. Then $MM^T=D+A(X)$,
  $M^TM=2I+A(L(X))$, and $\rk M=|V(X)|-c$, where $c$ is the number of bipartite connected
  components of $X$. Every eigenvalue of $L(X)$ is at least $-2$.
 \item \factlabel{[ST, Lemma 6.3]}{st:6.3}\ST{Lemma 6.3}
  (\path{SimpleGraph.lineGraph_connected_iff_of_no_isolated}, \path{SimpleGraph.degree_lineGraph}).
  If $X$ has at least one edge and no isolated vertices, then $X$ is connected if and only if
  $L(X)$ is connected, and the edge $\{u,w\}$ has degree $\deg(u)+\deg(w)-2$ in $L(X)$.
 \item \factlabel{[ST, Theorem 6.4]}{st:6.4}\ST{Theorem 6.4} (SP4;
  \path{SimpleGraph.regular_or_semiregular_of_lineGraph_regular}). If $X$ has at least one edge
  and no isolated vertices, and $L(X)$ is connected and $h$-regular, then $X$ is regular, or $X$
  is bipartite with constant degree on each part.
 \item \factlabel{[ST, Theorem 6.6]}{st:6.6}\ST{Theorem 6.6}
  (\path{Matrix.spectrumMultiset_fromBlocks_scalar}). Let $B$ be $P\times Q$ of rank $\rho$,
  let $\mu_1,\dots,\mu_\rho$ be the positive eigenvalues of $BB^T$ with repetition, and
  let $s,t\in\re$. Then the spectrum of $C=\pmat{sI&B\\B^T&tI}$ consists of
  $\gamma_j^{\pm}=\frac12\bigl(s+t\pm\sqrt{(s-t)^2+4\mu_j}\bigr)$ for $1\le j\le\rho$,
  of $s$ with multiplicity $|P|-\rho$ and of $t$ with multiplicity $|Q|-\rho$.
 \item \factlabel{[ST, Theorem 6.7]}{st:6.7}\ST{Theorem 6.7}
  (\path{SimpleGraph.lineGraph_cospectral_bipartite}). Let $X_1$ be a connected bipartite graph
  with an edge such that $L(X_1)$ is regular, and let $X_2$ be a graph without isolated vertices
  such that $L(X_2)$ is cospectral with $L(X_1)$. Then $X_2$ is connected, and precisely one of
  the following holds: $X_2$ is bipartite, and with the part sizes ordered $p\le q$ in each graph,
  the part sizes, the respective degrees and the spectra of $B_1B_1^T$ and $B_2B_2^T$ agree; or
  $X_1$ is not regular, $X_2$ is regular and not bipartite, and $|V(X_1)|=|V(X_2)|+1$.
 \item \factlabel{[ST, Corollary 6.8]}{st:6.8}\ST{Corollary 6.8}
  (\path{SimpleGraph.lineGraph_cospectral_bipartite_spectrum}). Let $X_1$ be a connected graph
  with a bipartition $(P,Q)$, biregular with degrees $(k_1,k_2)$, with $k_1>k_2>0$,
  $p=|P|\le q=|Q|$ and biadjacency matrix $B$ of rank $\rho$, and let $X_2$ be a connected
  regular graph that is not bipartite with $L(X_2)$ cospectral with $L(X_1)$. Then $k_1k_2$ is a
  positive eigenvalue of $BB^T$ of multiplicity 1, and if $\mu_2,\dots,\mu_\rho$ are the other
  positive eigenvalues of $BB^T$ with repetition, then $\Spec(X_2)$ consists of
  $\frac12(k_1+k_2)$, of $\pm\sqrt{\mu_j+\frac14(k_1-k_2)^2}$ for $2\le j\le\rho$, of
  $\frac12(k_1-k_2)$ with multiplicity $p-\rho$ and of $-\frac12(k_1-k_2)$ with multiplicity
  $q-\rho$.
 \item \factlabel{[ST, Theorem 7.7]}{st:7.7}\ST{Theorem 7.7}
  (\path{Matrix.gram_eq_of_two_eigenvalues}). Let $N$ be a $v\times b$ zero-one matrix with
  $v>1$, constant row sum $r>0$ and constant column sum $k>1$, and let $\lambda>0$ with
  $r(k-1)=\lambda(v-1)$. If $\Spec(NN^T)=\ms{(rk)^{1},(r-\lambda)^{v-1}}$, then
  $NN^T=(r-\lambda)I+\lambda J$. If moreover $r>\lambda$, then $b\ge v$ (Fisher's inequality,
  DS3).
 \item \factlabel{[ST, Theorem 7.10]}{st:7.10}\ST{Theorem 7.10}
  (\path{SimpleGraph.IsRegularOfDegree.aeval_eq_smul_allOnes_of_connected}). If $X$ is connected
  and $k$-regular on $n>0$ vertices, and the real polynomial $f$ vanishes at every eigenvalue
  of $X$ other than $k$, then $f(A(X))=\frac{f(k)}{n}J$.
\end{itemize}
\endgroup

\begingroup\raggedright
\medskip\noindent\textbf{Inputs from Mathlib and from the Hall--Connor section.}\label{facts:sec5}
\uses{hc:facts:finite}
We will use the following declarations of Mathlib at commit \path{082e2d37}
(4 September 2026), the version pinned in \path{lake-manifest.json}: \path{SimpleGraph.Iso.lineGraph} (an isomorphism of graphs induces one of
their line graphs), \path{SimpleGraph.lineGraph_adj_iff_exists} (two edges are adjacent in the
line graph exactly when they are distinct and share an endpoint),
\path{SimpleGraph.sum_degrees_eq_twice_card_edges}, \path{SimpleGraph.degree_compl},
\path{SimpleGraph.adjMatrix_mul_self_apply_self} (the diagonal of $A(X)^2$ is the degree),
\path{SimpleGraph.cycleGraph_adj'}, \path{SimpleGraph.Iso.completeGraph},
\path{SimpleGraph.completeBipartiteGraphCongr}, \path{Equiv.sumComm}, \path{Fintype.equivFin},
\path{Finset.card_le_univ}, \path{Finset.eq_of_subset_of_card_le},
\path{LinearEquiv.finrank_eq}, \path{Submodule.finrank_le}, \path{finrank_euclideanSpace_fin}
and the structure \path{Matrix.IsHadamard}. Double counting is the paragraph
\hyperref[hc:facts:finite]{(finite counting and order)} of the Hall--Connor section. The
design interface is Definition~\ref{hc:def:design} of that section, with its incidence matrix
$N$ of Definition~\ref{hc:def:matrices}. This is the one design interface of Section 5; the
toolkit's facts concern real matrices, and we apply them to the entrywise real cast of the
integer matrix $N$, which we also write $N$. The cast is taken in Lemma~\ref{lem:design},
assertion 3, and in Theorem~\ref{thm:5.8}, where \ST{Theorem 7.7} is applied to a real zero-one
matrix. The toolkit's Section 7 proposes the same names \path{Design} and
\path{Design.incMatrix} with a real incidence matrix; the shared declarations are those of the
Hall--Connor section, over $\ints$, and we use no toolkit result stated for the toolkit's
designs.\par
\endgroup

\begingroup\raggedright
\medskip\noindent\textbf{Inputs from Section 4.}\label{input:sec4}
Theorem 5.8 rests on Theorem 4.11. We recall what we will use from Section 4, where these
results are proved: Definition~\ref{def:CP}, Definition~\ref{def:represented} and
Theorem~\ref{thm:4.11}. Theorem~\ref{thm:5.8} and Lemma~\ref{lem:dense} depend on those labels
directly, and none of them depends on Section 5.
\begin{itemize}
 \item The \emph{cocktail party graph} $CP(m)$ (\path{SimpleGraph.cocktailParty}) has vertex set
  $\{1,\dots,m\}\times\{0,1\}$, and $(i,a)$ is adjacent to $(j,a')$ exactly when $i\ne j$. It is
  $(2m-2)$-regular on $2m$ vertices.
 \item A graph $X$ is \emph{represented in $\typ E_8$} (\path{SimpleGraph.IsRepresentedIn}) if
  there is a family $(\Zx_u)_{u\in V(X)}$ of vectors of $\Phi(\typ E_8)\sbs\re^8$ (model (ii) of
  Definition~\ref{def:3.2}) whose Gram matrix is $2I+A(X)$.
 \item Theorem 4.11 (Hoffman and Ray-Chaudhuri; \path{SimpleGraph.hoffman_rayChaudhuri}). If
  $X$ is connected and regular and $2I+A(X)$ is positive semidefinite, then $X$ is isomorphic to
  $L(Y)$ for some graph $Y$, or $X$ is isomorphic to $CP(m)$ for some $m$, or $X$ is represented
  in $\typ E_8$.
\end{itemize}
\endgroup

\begingroup\raggedright
\medskip\noindent\textbf{Inputs for Example 5.13.}\label{input:ex513}
The example needs one fact that this section does not prove, EX4 of
\path{cgss/prerequisites.md}, and the existence results EX1 to EX3 for its instances. We record
them as deferred inputs; none of them may depend on Section 5.
\begin{itemize}
 \item (EX4, deferred.) If $Y$ is a connected graph on more than three vertices that is not
  bipartite, and $Z$ is a bipartite graph without isolated vertices, then $L(Y)\not\cong L(Z)$.
  This is a case of Whitney's theorem on line graphs, which Mathlib does not have. The
  prerequisites prove it from four facts, which belong to Sections 3 and 4: the star closure of
  the standard representation of $L(Y)$ is $\typ D_n$ when $Y$ is connected, not bipartite and on
  $n$ vertices; that of $L(Z)$ is $\typ A_{m-1}$ when $Z$ is connected, bipartite and on $m$
  vertices; isomorphic graphs have isometric representations, so the star closures are
  isomorphic (LN3 and LA2 there), and $\typ D_n\not\cong\typ A_n$ for $n\ge4$; and $Z$ is
  connected, since $L(Z)\cong L(Y)$ is connected and $Z$ has no isolated vertices. \new\ The
  Lean proof of Example 5.13 avoids EX4, by a clique bound recorded at the step
  on non-isomorphism in the proof of Example~\ref{ex:5.13}; we keep EX4 for the route of the paper.
 \item (EX1 to EX3; computation C23.) Resolvable designs with parameters
  $(k(2k-1),k,4k^2-1,2k+1,1)$ exist for $k=2^d$ with $d\ge1$ and for $k=3$, and designs with
  parameters $(6,3,10,5,2)$ and $(28,21,36,27,20)$ exist. We will use these only in the remark after
  Example 5.13. Computation C23 gives in Lean the designs that the paper's instances $k=2$ and $k=4$
  need; their existence for every $k=2^d$ and for $k=3$ is a citation, outside the scope.
\end{itemize}
\endgroup

\subsection{Lemmas 5.1 to 5.7}

\medskip\noindent\textbf{Definitions and auxiliary lemmas.}

\begin{namedresult}{Definition}{(bipartitions and biadjacency matrices)}\label{def:bipartition}
 \lean{SimpleGraph.IsBiregularWith}\lean{SimpleGraph.biadjMatrix}\uses{}
 \assume{$X$ is a graph.}
 A \emph{bipartition} of $X$ is a pair $(P,Q)$ of disjoint sets with $P\cup Q=V(X)$ such that
 every edge has one end in $P$ and one in $Q$; in Mathlib, \path{X.IsBipartiteWith P Q}
 (\path{SimpleGraph.IsBipartiteWith}) with $Q=V(X)\diff P$. $X$ is bipartite if it has a bipartition. The bipartition is \emph{biregular
 with degrees} $(k_1,k_2)$ if every vertex of $P$ has degree $k_1$ and every vertex of $Q$ has
 degree $k_2$. The \emph{biadjacency matrix} of $X$ for $(P,Q)$ is the $P\times Q$ matrix $B$
 with $B(u,w)=A(X)(u,w)$.
\end{namedresult}

\begin{namedresult}{Definition}{(incidence graphs, flag graphs, characterization)}\label{def:flag}
 \lean{Design.incidenceGraph}\lean{Design.flagGraph}\lean{Design.FlagGraphCharacterized}
 \uses{hc:def:design}
 \leavevmode\par
 \assume{$D$ is a design with parameters $(v,k,b,r,\lambda)$.}
 Here $D$ is as in Definition~\ref{hc:def:design}, with points $P$ and block labels $\cB$. The
 \emph{incidence graph} $I(D)$ has vertex set $P\sqcup\cB$, and a point $p$ is adjacent to a
 block $B$ when $p\in B$; there are no other edges. The \emph{flag graph} is $F(D)=L(I(D))$. We
 say that $F(D)$ is \emph{characterized by its spectrum} if every graph cospectral with $F(D)$
 is isomorphic to $F(D_2)$ for some design $D_2$ with the parameters $(v,k,b,r,\lambda)$, on
 some finite types $P_2$ of points and $\cB_2$ of block labels with decidable equality. In Lean,
 $P_2$ and $\cB_2$ range over \texttt{Type}, as do the vertex types of all graphs of this section;
 the design $D_2$ of Theorem 5.8 lives on subtypes of the vertex type of a graph.
 \repair\ This is the sense of the paper's proofs and of its sources (errata, Implicit
 hypothesis 6); in the stronger sense ``determined up to isomorphism'' Corollaries 5.10 and 5.11
 and Remark 5.12 are false.
\end{namedresult}

\begin{namedresult}{Lemma}{(design identities)}\label{lem:design}
 \lean{Design.b_mul_k_eq_v_mul_r}\lean{Design.r_mul_k_sub_one_eq}\lean{Design.k_le_v}
 \lean{Design.one_le_lam}\lean{Design.lam_le_r}\lean{Design.lam_eq_r_iff}
 \lean{Design.incMatrix_mul_transpose}\lean{Design.spectrumMultiset_incMatrix_mul_transpose}
 \lean{Design.v_le_b_of_lam_lt_r}
 \uses{hc:def:design,hc:def:matrices,hc:facts:finite,facts:sec5,st:5.9,st:7.7}
 \assume{$D$ is a design with parameters $(v,k,b,r,\lambda)$ on $P$ and $\cB$, with incidence
  matrix $N$.}
 \prove{\begin{pfenum}
  \item $bk=vr$ and $r(k-1)=\lambda(v-1)$.
  \item $k\le v$, $\lambda\ge1$ and $r\ge\lambda$; moreover $r=\lambda$ if and only if $k=v$,
   if and only if every block contains every point.
  \item $NN^T=(r-\lambda)I+\lambda J$ and $\Spec(NN^T)=\ms{(rk)^{1},(r-\lambda)^{v-1}}$.
  \item If $r>\lambda$, then $b\ge v$.
 \end{pfenum}}
\end{namedresult}
\begin{proof}
 \pfsketch\ These are DS1 and DS3 of \path{cgss/prerequisites.md}: double counting, and the
 Fisher inequality of the toolkit's Theorem 7.7.
 \step{s5dsflags}{$bk=vr$.}
 \begin{proof}\pf\ Count the pairs $(p,B)$ with $p\in B$ by blocks and by points, with the
  double counting of \hyperref[hc:facts:finite]{(finite counting and order)}.~\qed\end{proof}
 \step{s5dspairs}{$r(k-1)=\lambda(v-1)$.}
 \begin{proof}\pf\ Fix a point $p$ and count the pairs $(q,B)$ with $q\ne p$ and $p,q\in B$, by
  blocks through $p$ and by points $q$, with the same double counting.~\qed\end{proof}
 \step{s5dskv}{$k\le v$.}
 \begin{proof}\pf\ A block is a subset of $P$ of size $k$; apply \path{Finset.card_le_univ}.~\qed\end{proof}
 \step{s5dslam}{$\lambda\ge1$.}
 \begin{proof}
  \step{s5dslamzero}{\assume{$\lambda=0$.}\prove{\textsc{False}.}}
  \begin{proof}\pf\ By \stepref{s5dspairs}, $r(k-1)=0$, and $k>1$ gives $r=0$. Then
   \stepref{s5dsflags} gives $bk=0$, against $b>0$ and $k>1$.~\qed\end{proof}
  \qedstep
  \begin{proof}\pf\ $\lambda$ is a natural number and \stepref{s5dslamzero}.~\qed\end{proof}
 \end{proof}
 \step{s5dsrlam}{$r\ge\lambda$, with equality if and only if $k=v$.}
 \begin{proof}\pf\ By \stepref{s5dspairs} and \stepref{s5dskv},
  $\lambda(v-1)=r(k-1)\le r(v-1)$, with equality exactly when $k=v$, since $r\ge1$ by
  \stepref{s5dspairs} and \stepref{s5dslam}; divide by $v-1>0$.~\qed\end{proof}
 \step{s5dsfull}{$k=v$ if and only if every block contains every point.}
 \begin{proof}\pf\ A block of size $k$ is a subset of $P$; apply
  \path{Finset.eq_of_subset_of_card_le} to it and $P$.~\qed\end{proof}
 \step{s5dsgram}{$NN^T=(r-\lambda)I+\lambda J$.}
 \begin{proof}\pf\ $(NN^T)(p,q)$ is the number of blocks containing $p$ and $q$, which is $r$
  for $p=q$ and $\lambda$ otherwise, by Definition~\ref{hc:def:design}.~\qed\end{proof}
 \step{s5dsspec}{$\Spec(NN^T)=\ms{(rk)^{1},(r-\lambda)^{v-1}}$.}
 \begin{proof}\pf\ \ST{Corollary 5.9} with $c=r-\lambda$, $t=\lambda$ and $n=v$ gives the
  eigenvalue $r-\lambda+\lambda v$ once, and $r-\lambda+\lambda v=rk$ by
  \stepref{s5dspairs}.~\qed\end{proof}
 \step{s5dsfisher}{\assume{$r>\lambda$.}\prove{$b\ge v$.}}
 \begin{proof}\pf\ \ST{Theorem 7.7} for the real zero-one $v\times b$ matrix $N$, whose row
  sums are $r>\lambda\ge1$ and column sums $k>1$ (Definition~\ref{hc:def:matrices}), with $v>1$,
  $\lambda>0$ by \stepref{s5dslam}, \stepref{s5dspairs}, \stepref{s5dsspec} and \asm{1}.~\qed\end{proof}
 \qedstep
 \begin{proof}\pf\ The preceding steps are the four assertions.~\qed\end{proof}
\end{proof}

\begin{namedresult}{Lemma}{(incidence and flag graphs)}\label{lem:incidence}
 \lean{Design.incidenceGraph_isBiregularWith}\lean{Design.incidenceGraph_connected}
 \lean{Design.incidenceGraph_isRegular_iff}\lean{Design.flagGraph_connected}
 \lean{Design.flagGraph_isRegularOfDegree}\lean{Design.card_flagGraph}
 \lean{Design.posSemidef_two_add_adjMatrix_flagGraph}\lean{Design.rank_two_add_adjMatrix_flagGraph}
 \uses{def:flag,def:bipartition,hc:def:matrices,lem:design,facts:sec5,st:3.5,st:6.2,st:6.3}
 \assume{$D$ is a design with parameters $(v,k,b,r,\lambda)$ on $P$ and $\cB$, with incidence
  matrix $N$.}
 \prove{\begin{pfenum}
  \item $(P,\cB)$ is a bipartition of $I(D)$, biregular with degrees $(r,k)$, whose
   biadjacency matrix is $N$; $I(D)$ has $vr$ edges and no isolated vertices.
  \item $I(D)$ is connected, and it is regular if and only if $r=k$.
  \item $F(D)$ is connected and $(r+k-2)$-regular on $vr$ vertices.
  \item $2I+A(F(D))$ is positive semidefinite of rank $v+b-1$.
 \end{pfenum}}
\end{namedresult}
\begin{proof}
 \pfsketch\ We will read the degrees off the definitions, and then apply the toolkit's facts on
 line graphs and incidence matrices.
 \step{s5igbip}{Assertion 1 holds.}
 \begin{proof}
  \step{s5igparts}{$(P,\cB)$ is a bipartition with biadjacency matrix $N$.}
  \begin{proof}\pf\ Every edge of $I(D)$ joins a point to a block, by Definition~\ref{def:flag},
   and $N(p,B)=1$ exactly when $p\in B$ (Definition~\ref{hc:def:matrices}).~\qed\end{proof}
  \step{s5igdeg}{Points have degree $r\ge1$ and blocks have degree $k\ge2$.}
  \begin{proof}\pf\ The neighbours of $p$ are the blocks through $p$, and those of $B$ are its
   $k$ points; $r\ge\lambda\ge1$ by Lemma~\ref{lem:design}, and $k>1$.~\qed\end{proof}
  \qedstep
  \begin{proof}\pf\ The number of edges is the number of pairs $(p,B)$ with $p\in B$, which is
   $vr$ by \stepref{s5igdeg} and double counting; there are no isolated vertices by
   \stepref{s5igdeg}.~\qed\end{proof}
 \end{proof}
 \step{s5igconn}{$I(D)$ is connected.}
 \begin{proof}
  \step{s5igpoints}{Any two points are joined by a path.}
  \begin{proof}\pf\ Two distinct points lie in $\lambda\ge1$ common blocks, by
   Lemma~\ref{lem:design}; such a block is a common neighbour.~\qed\end{proof}
  \qedstep
  \begin{proof}\pf\ $P$ is nonempty since $v>1$, and every block has a neighbouring point, since
   $k\ge2$; combine this with \stepref{s5igpoints}.~\qed\end{proof}
 \end{proof}
 \step{s5igreg}{$I(D)$ is regular if and only if $r=k$.}
 \begin{proof}\pf\ Both parts are nonempty ($v>1$, $b>0$), and the degrees are $r$ and $k$ by
  \stepref{s5igbip}.~\qed\end{proof}
 \step{s5igflag}{Assertion 3 holds.}
 \begin{proof}\pf\ Apply \ST{Lemma 6.3} to $I(D)$, which has an edge and no isolated vertices
  by \stepref{s5igbip} and is connected by \stepref{s5igconn}: $F(D)$ is connected, and an edge
  $\{p,B\}$ has degree $r+k-2$ in it. $F(D)$ has $|E(I(D))|=vr$ vertices.~\qed\end{proof}
 \step{s5igrank}{Assertion 4 holds.}
 \begin{proof}\pf\ Let $M$ be the incidence matrix of $I(D)$. By \ST{Theorem 6.2},
  $M^TM=2I+A(F(D))$ and $\rk M=(v+b)-1$, since $I(D)$ is one bipartite component by
  \stepref{s5igbip} and \stepref{s5igconn}. $M^TM$ is a Gram matrix, positive semidefinite
  with $\rk(M^TM)=\rk M$ by \ST{Theorem 3.5}.~\qed\end{proof}
 \qedstep
 \begin{proof}\pf\ Steps \stepref{s5igbip} to \stepref{s5igrank}.~\qed\end{proof}
\end{proof}

\begin{namedresult}{Lemma}{(regular graphs of large valency)}\label{lem:dense}
 \lean{SimpleGraph.eq_top_of_isRegularOfDegree_card_sub_one}
 \lean{SimpleGraph.nonempty_iso_cocktailParty}
 \uses{def:CP,facts:sec5}
 \assume{$X$ is a graph on $n$ vertices.}
 \prove{\begin{pfenum}
  \item If $n\ge1$ and $X$ is $(n-1)$-regular, then $X$ is complete, and $X\cong K_n$.
  \item If $n=2m$ with $m\ge1$ and $X$ is $(2m-2)$-regular, then $X\cong CP(m)$.
 \end{pfenum}}
\end{namedresult}
\begin{proof}
 \pfsketch\ \repair\ The paper does not state these, but its corollaries leave them to the
 reader (errata, Gaps 15 and 16). We will look at the non-neighbours of each vertex.
 \step{s5dncompl}{Every vertex has $n-1-\deg(u)$ non-neighbours other than itself.}
 \begin{proof}\pf\ \path{SimpleGraph.degree_compl}.~\qed\end{proof}
 \step{s5dncomplete}{Assertion 1 holds.}
 \begin{proof}\pf\ By \stepref{s5dncompl}, each vertex is adjacent to every other vertex, so
  $X$ is the complete graph on $V(X)$; apply \path{SimpleGraph.Iso.completeGraph} to
  \path{Fintype.equivFin}.~\qed\end{proof}
 \step{s5dncp}{Assertion 2 holds.}
 \begin{proof}
  \step{s5dnpartner}{Each vertex $u$ has a unique non-neighbour $u^*\ne u$, and $u^{**}=u$.}
  \begin{proof}\pf\ \stepref{s5dncompl} gives exactly one, and non-adjacency is
   symmetric.~\qed\end{proof}
  \step{s5dnenum}{\pick\ a bijection $e:V(X)\to\{1,\dots,m\}\times\{0,1\}$ such that $e(u)$ and
   $e(u^*)$ have the same first coordinate for every $u$.}
  \begin{proof}\pf\ The pairs $\{u,u^*\}$ partition $V(X)$ into $m$ pairs, by
   \stepref{s5dnpartner}; number the pairs with \path{Fintype.equivFin} and the two members of
   each pair by $0$ and $1$.~\qed\end{proof}
  \qedstep
  \begin{proof}\pf\ For $u\ne w$, $u$ and $w$ are adjacent if and only if $w\ne u^*$, by
   \stepref{s5dnpartner}, that is, when $e(u)$ and $e(w)$ have different first coordinates by
   \stepref{s5dnenum}; this is adjacency in $CP(m)$.~\qed\end{proof}
 \end{proof}
 \qedstep
 \begin{proof}\pf\ Steps \stepref{s5dncomplete} and \stepref{s5dncp}.~\qed\end{proof}
\end{proof}

\begin{namedresult}{Lemma}{(connected $2$-regular graphs)}\label{lem:cycle}
 \lean{SimpleGraph.nonempty_iso_cycleGraph}\uses{facts:sec5}
 \assume{$X$ is a connected $2$-regular graph on $n$ vertices.}
 \prove{$n\ge3$ and $X\cong C_n$, the cycle \path{SimpleGraph.cycleGraph n}.}
\end{namedresult}
\begin{proof}
 \pfsketch\ \repair\ The paper leaves this to the reader in Corollaries 5.10 and 5.11 (errata,
 Gap 16). We will walk around $X$ without turning back.
 \step{s5cyseq}{\pick\ adjacent vertices $u_0,u_1$, and for $i\ge1$ let $u_{i+1}$ be the
  neighbour of $u_i$ other than $u_{i-1}$.}
 \begin{proof}\pf\ $X$ is nonempty since it is connected, and a vertex has two neighbours; each
  $u_i$ with $i\ge1$ is adjacent to $u_{i-1}$, so it has exactly one other
  neighbour.~\qed\end{proof}
 \step{s5cyfirst}{\pick\ the least $j$ such that $u_j=u_i$ for some $i<j$. Then $i=0$ and
  $j\ge3$.}
 \begin{proof}
  \step{s5cyexists}{Such a $j$ exists, with $j\le n$.}
  \begin{proof}\pf\ The vertices $u_0,\dots,u_n$ are $n+1$ vertices of a set of size $n$.~\qed\end{proof}
  \step{s5cynotlater}{\assume{$u_j=u_i$ with $1\le i<j$.}\prove{\textsc{False}.}}
  \begin{proof}\pf\ By \stepref{s5cyseq}, $u_{i+1}\ne u_i$ and $u_{i+2}\ne u_i$, so $j\ge i+3$;
   by the minimality of $j$, $u_{i-1}$, $u_{i+1}$ and $u_{j-1}$ are then three distinct
   neighbours of $u_i=u_j$, against $2$-regularity.~\qed\end{proof}
  \qedstep
  \begin{proof}\pf\ Steps \stepref{s5cyexists} and \stepref{s5cynotlater} give $i=0$; and
   $u_1\ne u_0$, $u_2\ne u_0$ by \stepref{s5cyseq}, so $j\ge3$.~\qed\end{proof}
 \end{proof}
 \step{s5cyclosed}{The set $W=\{u_0,\dots,u_{j-1}\}$ contains all neighbours of its members.}
 \begin{proof}\pf\ For $1\le i\le j-1$ the two neighbours of $u_i$ are $u_{i-1}$ and $u_{i+1}$,
  with $u_j=u_0$; the two neighbours of $u_0$ are $u_1$ and $u_{j-1}$, distinct since $j\ge3$,
  by \stepref{s5cyfirst}.~\qed\end{proof}
 \qedstep
 \begin{proof}\pf\ By \stepref{s5cyclosed} and connectedness, $W=V(X)$, so $j=n\ge3$. By the
  same step, $u_i\mapsto i$ is a bijection $V(X)\to\{0,\dots,n-1\}$ under which $u_i$ and $u_{i'}$
  are adjacent exactly when $i-i'\equiv\pm1\pmod n$; this is adjacency in
  $C_n$ by \path{SimpleGraph.cycleGraph_adj'}.~\qed\end{proof}
\end{proof}

\begin{namedresult}{Lemma}{(isolated vertices)}\label{lem:isolated}
 \lean{SimpleGraph.exists_forall_degree_pos_nonempty_iso_lineGraph}\uses{facts:sec5}
 \assume{$Y$ is a graph.}
 \prove{There is a graph $Y'$ without isolated vertices such that $L(Y')\cong L(Y)$.}
\end{namedresult}
\begin{proof}
 \pfsketch\ \repair\ Isolated vertices are invisible to the line graph (errata, Implicit
 hypotheses 3 and 5). We will delete them.
 \pflet{$W$ be the set of vertices of $Y$ of positive degree, and $Y'$ the subgraph of $Y$
  induced on $W$.}
 \step{s5isnone}{$Y'$ has no isolated vertices.}
 \begin{proof}\pf\ A neighbour of $w\in W$ has positive degree, so it lies in $W$.~\qed\end{proof}
 \step{s5isedges}{Every edge of $Y$ has both ends in $W$, so $\{u,w\}\mapsto\{u,w\}$ is a
  bijection $E(Y')\to E(Y)$.}
 \begin{proof}\pf\ The ends of an edge have positive degree; $Y'$ is induced.~\qed\end{proof}
 \qedstep
 \begin{proof}\pf\ The bijection of \stepref{s5isedges} preserves the ends of edges, so it
  preserves adjacency in the line graphs by \path{SimpleGraph.lineGraph_adj_iff_exists}; with
  \stepref{s5isnone} this is the claim.~\qed\end{proof}
\end{proof}

\medskip\noindent\textbf{Lemmas 5.1 to 5.5.}

\begin{cgss}{Lemma}{5.1}\label{lem:5.1}
 \lean{SimpleGraph.IsRegularOfDegree.largestEigenvalue_eq}
 \lean{SimpleGraph.IsRegularOfDegree.multiplicity_eq_card_components}\uses{st:5.3}
 \assume{$X$ is a $k$-regular graph on a nonempty vertex set.}
 \prove{$\lmax(X)=k$, and $\mult_X(k)$ is the number of connected components of $X$.}
\end{cgss}
\begin{proof}
 \pfsketch\ \repair\ The paper quotes this as well known (errata, Gap 8); the toolkit proves it
 through the Laplacian $kI-A(X)$.
 \qedstep
 \begin{proof}\pf\ \ST{Theorem 5.3}.~\qed\end{proof}
\end{proof}

\begin{cgss}{Lemma}{5.2}\label{lem:5.2}
 \lean{SimpleGraph.isRegular_iff_sum_sq_eigenvalues}\uses{st:5.4}
 \assume{$X$ is a graph on $n\ge1$ vertices.}
 \prove{$X$ is regular if and only if $\sum_{\mu\in\Spec(X)}\mu^2=n\lmax(X)$, where the sum
  counts multiplicities.}
\end{cgss}
\begin{proof}
 \pfsketch\ \repair\ The paper quotes this as well known (errata, Gap 8); the toolkit proves it
 from the Rayleigh quotient of $\Zj$.
 \qedstep
 \begin{proof}\pf\ \ST{Theorem 5.4}.~\qed\end{proof}
\end{proof}

\begin{cgss}{Lemma}{5.3}\label{lem:5.3}
 \lean{SimpleGraph.IsRegularOfDegree.isBipartite_iff_neg_eigenvalue}\uses{st:5.6}
 \assume{$X$ is a connected $k$-regular graph.}
 \prove{$-k\in\Spec(X)$ if and only if $X$ is bipartite, and then $\mult_X(-k)=1$.}
\end{cgss}
\begin{proof}
 \pfsketch\ \repair\ The paper quotes this as well known (errata, Gap 8); the toolkit proves it
 through the signless Laplacian $kI+A(X)$.
 \qedstep
 \begin{proof}\pf\ \ST{Theorem 5.6}.~\qed\end{proof}
\end{proof}

\begin{cgss}{Lemma}{5.4}\label{lem:5.4}
 \lean{SimpleGraph.regular_or_isBiregularWith_of_lineGraph_regular}\uses{def:bipartition,st:6.4}
 \assume{\begin{pfenum}
  \item $X$ is a graph without isolated vertices.
  \item $L(X)$ is connected and regular.
 \end{pfenum}}
 \prove{$X$ is regular, or $X$ has a bipartition that is biregular with degrees $(k_1,k_2)$
  for some $k_1\ne k_2$.}
\end{cgss}
\begin{proof}
 \pfsketch\ \repair\ The paper quotes this as well known (errata, Gap 8), and it needs
 \asm[1]{0} (errata, Implicit hypothesis 3: $K_2$ with an isolated vertex is a counterexample).
 The toolkit proves it by counting, as its Theorem 6.4, which also assumes an edge and asserts
 $k_1+k_2=h+2$ and $pk_1=qk_2=|E(X)|$ for the part sizes $p$, $q$. We will derive the edge and
 $k_1\ne k_2$; the form without an edge is the one that Theorem 5.8 uses, so this lemma has its
 own Lean name.
 \step{s554edge}{$X$ has an edge.}
 \begin{proof}\pf\ A connected graph is nonempty, and the vertices of $L(X)$ are the edges of
  $X$.~\qed\end{proof}
 \step{s554reg}{\case{$X$ is regular.}}
 \begin{proof}\pf\ This is the first alternative.~\qed\end{proof}
 \step{s554notreg}{\case{$X$ is not regular.}}
 \begin{proof}\pf\ By \ST{Theorem 6.4}, with \stepref{s554edge}, \asm[1]{0} and \asm[2]{0}, $X$
  has a bipartition that is biregular with some degrees $(k_1,k_2)$, and $k_1\ne k_2$ since $X$
  is not regular.~\qed\end{proof}
 \qedstep
 \begin{proof}\pf\ The cases \stepref{s554reg} and \stepref{s554notreg} are
  exhaustive.~\qed\end{proof}
\end{proof}

\begin{cgss}{Lemma}{5.5}\label{lem:5.5}
 \lean{SimpleGraph.cospectral_of_lineGraph_cospectral}
 \uses{lem:5.1,facts:sec5,st:2.2,st:2.4,st:2.6,st:3.1,st:6.1,st:6.2,st:6.3}
 \assume{\begin{pfenum}
  \item $X$ and $X'$ are regular graphs, each with at least one edge.
  \item $L(X)$ and $L(X')$ are cospectral.
 \end{pfenum}}
 \prove{$X$ and $X'$ are cospectral.}
\end{cgss}
\begin{proof}
 \pfsketch\ \repair\ The paper states that the spectrum of a regular graph is determined by that
 of its line graph; this needs an edge, since all edgeless graphs have the same line graph
 (errata, Implicit hypothesis 3). We will follow the paper: the largest eigenvalue of the line
 graph gives the valency, its order gives the number of vertices, and $MM^T$ and $M^TM$ have the
 same nonzero eigenvalues.
 \pflet{$k$, $k'$ be the valencies, $n$, $n'$ the numbers of vertices and $m$, $m'$ the numbers
  of edges of $X$, $X'$, and $M$, $M'$ their incidence matrices.}
 \step{s555m}{$m=m'$.}
 \begin{proof}\pf\ The cospectral graphs $L(X)$, $L(X')$ have the same number of vertices, by
  \ST{Proposition 2.6}.~\qed\end{proof}
 \step{s555lk}{$L(X)$ is $(2k-2)$-regular and $L(X')$ is $(2k'-2)$-regular, with $k,k'\ge1$.}
 \begin{proof}\pf\ An edge gives $k,k'\ge1$, so there are no isolated vertices; apply the
  degree formula of \ST{Lemma 6.3}.~\qed\end{proof}
 \step{s555k}{$k=k'$.}
 \begin{proof}\pf\ By Lemma~\ref{lem:5.1} and \stepref{s555lk}, $\lmax(L(X))=2k-2$ and
  $\lmax(L(X'))=2k'-2$; cospectral graphs have the same largest eigenvalue by
  \ST{Proposition 2.6} and \ST{Lemma 2.4}.~\qed\end{proof}
 \step{s555n}{$n=n'$.}
 \begin{proof}\pf\ $nk=2m$ and $n'k'=2m'$ by \path{SimpleGraph.sum_degrees_eq_twice_card_edges};
  use \stepref{s555m}, \stepref{s555k} and $k\ge1$.~\qed\end{proof}
 \step{s555gram}{$MM^T=kI+A(X)$, $M^TM=2I+A(L(X))$, $M'M'^T=k'I+A(X')$ and
  $M'^TM'=2I+A(L(X'))$.}
 \begin{proof}\pf\ \ST{Theorem 6.2} for $X$ and for $X'$, with the degree matrices $kI$ and
  $k'I$.~\qed\end{proof}
 \step{s555nonzero}{$\mult_{kI+A(X)}(\mu)=\mult_{kI+A(X')}(\mu)$ for every $\mu\ne0$.}
 \begin{proof}
  \step{s555nzX}{$\mult_{kI+A(X)}(\mu)=\mult_{2I+A(L(X))}(\mu)$ and
   $\mult_{k'I+A(X')}(\mu)=\mult_{2I+A(L(X'))}(\mu)$.}
  \begin{proof}\pf\ \ST{Theorem 6.1} for $M$, $M^T$ and for $M'$, $M'^T$, with
   \stepref{s555gram}, since $\mu\ne0$.~\qed\end{proof}
  \step{s555nzL}{$\mult_{2I+A(L(X))}(\mu)=\mult_{2I+A(L(X'))}(\mu)$.}
  \begin{proof}\pf\ Both equal the multiplicity of $\mu-2$ in the common spectrum of $L(X)$
   and $L(X')$, by \ST{Theorem 3.1} and \ST{Proposition 2.6}.~\qed\end{proof}
  \qedstep
  \begin{proof}\pf\ Steps \stepref{s555nzX} and \stepref{s555nzL}, with $k=k'$ from
   \stepref{s555k}.~\qed\end{proof}
 \end{proof}
 \step{s555zero}{$\mult_{kI+A(X)}(0)=\mult_{kI+A(X')}(0)$.}
 \begin{proof}\pf\ By \ST{Theorem 2.2}, each multiplicity at $0$ is the order $n=n'$ minus the
  sum of the multiplicities at the nonzero eigenvalues; use \stepref{s555n} and
  \stepref{s555nonzero}.~\qed\end{proof}
 \qedstep
 \begin{proof}\pf\ By \stepref{s555nonzero}, \stepref{s555zero} and \ST{Proposition 2.6},
  $\Spec(kI+A(X))=\Spec(kI+A(X'))$. Shifting by $-k$ with \ST{Theorem 3.1} gives
  $\Spec(X)=\Spec(X')$, and \ST{Proposition 2.6} gives cospectrality.~\qed\end{proof}
\end{proof}

\medskip\noindent\textbf{Lemma 5.6 and Remark 5.7.}

We will need the paper's analysis of a connected bipartite graph with a regular line graph
twice, for $X_1$ and for $X_2$ in Lemma 5.6. We state it as two lemmas, so that no step is cited
from inside another proof.

\begin{namedresult}{Lemma}{(a bipartite root graph)}\label{lem:rootgraph}
 \lean{SimpleGraph.IsBiregularWith.lineGraph_isRegularOfDegree}
 \lean{SimpleGraph.IsBiregularWith.signlessLapMatrix_eq_fromBlocks}
 \lean{SimpleGraph.IsBiregularWith.multiplicity_zero_signlessLapMatrix}
 \lean{SimpleGraph.IsBiregularWith.card_eq_rank_add_one}
 \lean{SimpleGraph.IsBiregularWith.spectrumMultiset_map_quadratic_signlessLapMatrix}
 \uses{def:bipartition,lem:5.1,facts:sec5,hc:facts:finite,st:2.2,st:3.1,st:3.2,st:3.5,st:6.1,st:6.2,st:6.3,st:6.6}
 \assume{\begin{pfenum}
  \item $Z$ is a connected graph with at least one edge.
  \item $(P,Q)$ is a bipartition of $Z$, biregular with degrees $(k_1,k_2)$, with biadjacency
   matrix $B$.
  \item $M$ is the incidence matrix of $Z$ and $C=MM^T$.
 \end{pfenum}}
 \prove{\begin{pfenum}
  \item $Z$ has no isolated vertices, $k_1,k_2\ge1$, and $|P|k_1=|Q|k_2=|E(Z)|$.
  \item $L(Z)$ is connected and $(k_1+k_2-2)$-regular, and $\lmax(L(Z))=k_1+k_2-2$.
  \item In block form for $V(Z)=P\cup Q$, $C=\pmat{k_1I&B\\B^T&k_2I}$, and
   $BB^T\Zj=k_1k_2\Zj$.
  \item $\mult_C(0)=1$.
  \item $|V(Z)|=\rk(2I+A(L(Z)))+1$.
  \item Let $\rho=\rk B$, let $\Lambda$ be the multiset of positive eigenvalues of $BB^T$, and
   let $f(z)=(z-k_1)(z-k_2)$. Then $\Lambda$ has $\rho$ members, and the image of $\Spec(C)$
   under $f$ is $\Lambda+\Lambda+\ms{0^{|V(Z)|-2\rho}}$.
 \end{pfenum}}
\end{namedresult}
\begin{proof}
 \pfsketch\ This is the first part of the paper's proof of Lemma 5.6. The kernel of $C$ is the
 image of the eigenspace of $C$ for $k_1+k_2$ under an invertible map, and $k_1+k_2$ is a simple
 eigenvalue since $L(Z)$ is connected.
 \step{s5rgiso}{Assertion 1 holds.}
 \begin{proof}
  \step{s5rgnoiso}{$Z$ has no isolated vertices.}
  \begin{proof}\pf\ An edge gives two vertices, so every vertex $u$ has a vertex $w\ne u$ and a
   walk from $u$ to $w$, by \asm[1]{0}; its first edge is at $u$.~\qed\end{proof}
  \step{s5rgdeg}{$P$ and $Q$ are nonempty and $k_1,k_2\ge1$.}
  \begin{proof}\pf\ An edge has an end in $P$ and one in $Q$ by \asm[2]{0}; the degrees of
   these ends are positive by \stepref{s5rgnoiso}.~\qed\end{proof}
  \qedstep
  \begin{proof}\pf\ Each edge has exactly one end in $P$ and one in $Q$, so double counting the
   pairs $(u,e)$ with $u\in P$ an end of $e$, and the pairs $(w,e)$ with $w\in Q$ an end of $e$,
   gives $|P|k_1=|E(Z)|$ and $|Q|k_2=|E(Z)|$, by
   \hyperref[hc:facts:finite]{(finite counting and order)}.~\qed\end{proof}
 \end{proof}
 \step{s5rgline}{Assertion 2 holds.}
 \begin{proof}\pf\ \ST{Lemma 6.3} applies by \asm[1]{0} and \stepref{s5rgiso}; an edge has one
  end in $P$ and one in $Q$, so its degree in $L(Z)$ is $k_1+k_2-2$. Then
  Lemma~\ref{lem:5.1} gives $\lmax$.~\qed\end{proof}
 \step{s5rgblock}{Assertion 3 holds.}
 \begin{proof}
  \step{s5rgblockC}{$C=\pmat{k_1I&B\\B^T&k_2I}$.}
  \begin{proof}\pf\ By \ST{Theorem 6.2}, $C=D+A(Z)$ with $D$ the degree matrix, which is $k_1I$
   on $P$ and $k_2I$ on $Q$ by \asm[2]{0}; $A(Z)$ vanishes on $P\times P$ and $Q\times Q$ by
   \asm[2]{0}, and its $P\times Q$ block is $B$ by Definition~\ref{def:bipartition}.~\qed\end{proof}
  \step{s5rgblockj}{$B\Zj=k_1\Zj$ and $B^T\Zj=k_2\Zj$.}
  \begin{proof}\pf\ The row sum of $B$ at $u\in P$ is $\deg(u)=k_1$, and the column sum at $w\in
   Q$ is $\deg(w)=k_2$, by \asm[2]{0}.~\qed\end{proof}
  \qedstep
  \begin{proof}\pf\ \stepref{s5rgblockC}, and $BB^T\Zj=k_2B\Zj=k_1k_2\Zj$ by
   \stepref{s5rgblockj}.~\qed\end{proof}
 \end{proof}
 \step{s5rgsimple}{$\mult_C(k_1+k_2)=1$.}
 \begin{proof}
  \step{s5rgsimpleA}{$\mult_C(k_1+k_2)=\mult_{2I+A(L(Z))}(k_1+k_2)$.}
  \begin{proof}\pf\ \ST{Theorem 6.1} for $M$ and $M^T$, since $k_1+k_2\ne0$, with
   $M^TM=2I+A(L(Z))$ from \ST{Theorem 6.2}.~\qed\end{proof}
  \step{s5rgsimpleB}{$\mult_{2I+A(L(Z))}(k_1+k_2)=\mult_{L(Z)}(k_1+k_2-2)=1$.}
  \begin{proof}\pf\ The first equality is \ST{Theorem 3.1}; the second is Lemma~\ref{lem:5.1},
   since $L(Z)$ is connected and $(k_1+k_2-2)$-regular by \stepref{s5rgline}.~\qed\end{proof}
  \qedstep
  \begin{proof}\pf\ Steps \stepref{s5rgsimpleA} and \stepref{s5rgsimpleB}.~\qed\end{proof}
 \end{proof}
 \step{s5rgmap}{The map $\phi(\Zx,\Zy)=(-k_2\Zx,k_1\Zy)$ on $\re^P\oplus\re^Q$ is invertible and
  carries $\ker(C-(k_1+k_2)I)$ onto $\ker C$.}
 \begin{proof}
  \step{s5rgeigen}{$C(\Zx,\Zy)=(k_1+k_2)(\Zx,\Zy)$ if and only if $B\Zy=k_2\Zx$ and
   $B^T\Zx=k_1\Zy$.}
  \begin{proof}\pf\ Block multiplication with \stepref{s5rgblock}: $C(\Zx,\Zy)=(k_1\Zx+B\Zy,
   B^T\Zx+k_2\Zy)$.~\qed\end{proof}
  \step{s5rgkernel}{$C\phi(\Zx,\Zy)=0$ if and only if $B\Zy=k_2\Zx$ and $B^T\Zx=k_1\Zy$.}
  \begin{proof}\pf\ Block multiplication gives
   $C\phi(\Zx,\Zy)=\bigl(k_1(B\Zy-k_2\Zx),\,k_2(k_1\Zy-B^T\Zx)\bigr)$, and $k_1,k_2\ge1$ by
   \stepref{s5rgiso}.~\qed\end{proof}
  \qedstep
  \begin{proof}\pf\ $\phi$ is invertible since $k_1,k_2\ne0$; use \stepref{s5rgeigen} and
   \stepref{s5rgkernel}.~\qed\end{proof}
 \end{proof}
 \step{s5rgzero}{Assertion 4 holds.}
 \begin{proof}\pf\ By \ST{Theorem 2.2}, $\mult_C(0)=\dim\ker C$ and
  $\mult_C(k_1+k_2)=\dim\ker(C-(k_1+k_2)I)$; these are equal by \stepref{s5rgmap} and
  \path{LinearEquiv.finrank_eq}, and the second is 1 by \stepref{s5rgsimple}.~\qed\end{proof}
 \step{s5rgrank}{Assertion 5 holds.}
 \begin{proof}\pf\ $\rk C=|V(Z)|-1$ by \ST{Theorem 3.2} and \stepref{s5rgzero}, and
  $\rk C=\rk(MM^T)=\rk(M^TM)=\rk(2I+A(L(Z)))$ by \ST{Theorem 3.5} and
  \ST{Theorem 6.2}.~\qed\end{proof}
 \step{s5rgquad}{Assertion 6 holds.}
 \begin{proof}
  \pflet{$\mu_1,\dots,\mu_\rho$ be the members of $\Lambda$, which has $\rho$ members by
   \ST{Theorem 6.1}, and
   $\gamma_j^{\pm}=\frac12\bigl(k_1+k_2\pm\sqrt{(k_1-k_2)^2+4\mu_j}\bigr)$.}
  \step{s5rgroots}{$f(\gamma_j^+)=f(\gamma_j^-)=\mu_j$ for $1\le j\le\rho$.}
  \begin{proof}\pf\ $\gamma_j^{\pm}$ are the roots of $z^2-(k_1+k_2)z+k_1k_2-\mu_j$, by the
   quadratic formula, and $f(z)-\mu_j$ is this polynomial.~\qed\end{proof}
  \step{s5rgspecC}{$\Spec(C)$ consists of the $\gamma_j^{\pm}$, of $k_1$ with multiplicity
   $|P|-\rho$ and of $k_2$ with multiplicity $|Q|-\rho$.}
  \begin{proof}\pf\ \ST{Theorem 6.6} with $s=k_1$, $t=k_2$, and \stepref{s5rgblock}.~\qed\end{proof}
  \qedstep
  \begin{proof}\pf\ Apply $f$ to \stepref{s5rgspecC}: by \stepref{s5rgroots} and
   $f(k_1)=f(k_2)=0$, the image is $\Lambda+\Lambda+\ms{0^{|P|+|Q|-2\rho}}$, and
   $|P|+|Q|=|V(Z)|$.~\qed\end{proof}
 \end{proof}
 \qedstep
 \begin{proof}\pf\ Steps \stepref{s5rgiso}, \stepref{s5rgline}, \stepref{s5rgblock},
  \stepref{s5rgzero}, \stepref{s5rgrank} and \stepref{s5rgquad}.~\qed\end{proof}
\end{proof}

\begin{namedresult}{Lemma}{(two bipartite root graphs)}\label{lem:tworoot}
 \lean{SimpleGraph.IsBiregularWith.eq_of_lineGraph_cospectral}
 \uses{lem:rootgraph,facts:sec5,st:2.2,st:2.4,st:2.6,st:3.1,st:3.2,st:6.1,st:6.2}
 \assume{\begin{pfenum}
  \item $Z$ and $Z'$ are connected graphs, each with at least one edge.
  \item $(P,Q)$ is a bipartition of $Z$, biregular with degrees $(k_1,k_2)$, with biadjacency
   matrix $B$, and $|P|\le|Q|$.
  \item $(P',Q')$ is a bipartition of $Z'$, biregular with degrees $(k_1',k_2')$, with
   biadjacency matrix $B'$, and $|P'|\le|Q'|$.
  \item $L(Z)$ and $L(Z')$ are cospectral.
 \end{pfenum}}
 \prove{$|P|=|P'|$, $|Q|=|Q'|$, $k_1=k_1'$, $k_2=k_2'$, and $\Spec(BB^T)=\Spec(B'B'^T)$.}
\end{namedresult}
\begin{proof}
 \pfsketch\ This is the paper's claim that the spectrum of $L(X_1)$ determines $p$, $q$, $k_1$,
 $k_2$ and $\Spec(B_1B_1^T)$. We will recover the sizes and degrees from the number of
 edges, the valency and the rank, and then read $\Spec(BB^T)$ off $\Spec(C)$ through the
 quadratic $f(z)=(z-k_1)(z-k_2)$.
 \pflet{$m=|E(Z)|$, $N=|V(Z)|$, $C=MM^T$ for the incidence matrix $M$ of $Z$, and
  $m'=|E(Z')|$, $N'=|V(Z')|$, $C'=M'M'^T$ for the incidence matrix $M'$ of $Z'$.}
 \step{s5trm}{$m=m'$.}
 \begin{proof}\pf\ These are the numbers of vertices of the cospectral graphs $L(Z)$, $L(Z')$;
  apply \ST{Proposition 2.6}.~\qed\end{proof}
 \step{s5trs}{$k_1+k_2=k_1'+k_2'$.}
 \begin{proof}\pf\ $k_1+k_2-2$ and $k_1'+k_2'-2$ are the largest eigenvalues of $L(Z)$ and
  $L(Z')$, by Lemma~\ref{lem:rootgraph}; cospectral graphs have the same largest eigenvalue by
  \ST{Proposition 2.6} and \ST{Lemma 2.4}.~\qed\end{proof}
 \step{s5trN}{$N=N'$.}
 \begin{proof}
  \step{s5trrank}{$\rk(2I+A(L(Z)))=\rk(2I+A(L(Z')))$.}
  \begin{proof}\pf\ \ST{Theorem 3.2} with $f(z)=z+2$ expresses each rank through the common
   multiplicities of the two line graphs, which agree by \ST{Proposition 2.6}.~\qed\end{proof}
  \qedstep
  \begin{proof}\pf\ Lemma~\ref{lem:rootgraph}, assertion 5, for $Z$ and $Z'$, and
   \stepref{s5trrank}.~\qed\end{proof}
 \end{proof}
 \step{s5trprod}{$k_1k_2N=m(k_1+k_2)$ and $k_1'k_2'N'=m'(k_1'+k_2')$.}
 \begin{proof}\pf\ By Lemma~\ref{lem:rootgraph}, assertion 1, for $Z$ and for $Z'$,
  $N=|P|+|Q|=m/k_1+m/k_2$ and $N'=m'/k_1'+m'/k_2'$; multiply by $k_1k_2$ and by
  $k_1'k_2'$.~\qed\end{proof}
 \step{s5trorder}{$k_1\ge k_2$ and $k_1'\ge k_2'$.}
 \begin{proof}\pf\ By Lemma~\ref{lem:rootgraph}, assertion 1, for $Z$ and for $Z'$,
  $|P|k_1=|Q|k_2=m>0$ and $|P'|k_1'=|Q'|k_2'=m'>0$, with $|P|\le|Q|$ and
  $|P'|\le|Q'|$.~\qed\end{proof}
 \step{s5trdeg}{$k_1=k_1'$ and $k_2=k_2'$.}
 \begin{proof}\pf\ By \stepref{s5trm} to \stepref{s5trprod}, the pairs have the same sum and the
  same product, so $(k_1-k_2)^2=(k_1'-k_2')^2$; both differences are nonnegative by
  \stepref{s5trorder}.~\qed\end{proof}
 \step{s5trsize}{$|P|=|P'|$ and $|Q|=|Q'|$.}
 \begin{proof}\pf\ $|P|=m/k_1$, $|Q|=m/k_2$, $|P'|=m'/k_1'$ and $|Q'|=m'/k_2'$ by
  Lemma~\ref{lem:rootgraph}, assertion 1, for $Z$ and $Z'$; use \stepref{s5trm} and
  \stepref{s5trdeg}.~\qed\end{proof}
 \step{s5trC}{$\Spec(C)=\Spec(C')$.}
 \begin{proof}
  \step{s5trCnz}{$\mult_C(\gamma)=\mult_{C'}(\gamma)$ for every $\gamma\ne0$.}
  \begin{proof}\pf\ By \ST{Theorem 6.1} and \ST{Theorem 6.2}, $\mult_C(\gamma)=
   \mult_{2I+A(L(Z))}(\gamma)$ and $\mult_{C'}(\gamma)=\mult_{2I+A(L(Z'))}(\gamma)$; these agree
   by \ST{Theorem 3.1} and \ST{Proposition 2.6}.~\qed\end{proof}
  \qedstep
  \begin{proof}\pf\ $\mult_C(0)=\mult_{C'}(0)=1$ by Lemma~\ref{lem:rootgraph}, assertion 4; with
   \stepref{s5trCnz} and \ST{Proposition 2.6}.~\qed\end{proof}
 \end{proof}
 \pflet{$f(z)=(z-k_1)(z-k_2)$, and let $\Lambda$ and $\Lambda'$ be the multisets of positive
  eigenvalues of $BB^T$ and $B'B'^T$, and $\rho=\rk B$ and $\rho'=\rk B'$.}
 \step{s5trmap}{The images of $\Spec(C)$ and $\Spec(C')$ under $f$ are
  $\Lambda+\Lambda+\ms{0^{N-2\rho}}$ and $\Lambda'+\Lambda'+\ms{0^{N-2\rho'}}$.}
 \begin{proof}\pf\ Lemma~\ref{lem:rootgraph}, assertion 6, for $Z$, and for $Z'$, whose
  polynomial $(z-k_1')(z-k_2')$ is $f$ by \stepref{s5trdeg}, with $N'=N$ by
  \stepref{s5trN}.~\qed\end{proof}
 \step{s5trLambda}{$\Lambda=\Lambda'$.}
 \begin{proof}\pf\ By \stepref{s5trC} the two images agree. For $\mu>0$, the count of $\mu$ in
  them is $2\,\mathrm{count}_\Lambda(\mu)$ and $2\,\mathrm{count}_{\Lambda'}(\mu)$, by
  \stepref{s5trmap}; and $\Lambda$, $\Lambda'$ have only positive
  members.~\qed\end{proof}
 \qedstep
 \begin{proof}\pf\ By \ST{Theorem 6.1}, $\Spec(BB^T)=\Lambda+\ms{0^{|P|-\rho}}$ and
  $\Spec(B'B'^T)=\Lambda'+\ms{0^{|P'|-\rho'}}$; these agree by \stepref{s5trLambda}, which gives
  $\rho=\rho'$, and \stepref{s5trsize}. With \stepref{s5trdeg} and \stepref{s5trsize} this is
  the claim.~\qed\end{proof}
\end{proof}

\begin{namedresult}{Lemma}{(biregular bipartitions)}\label{lem:biregular}
 \lean{SimpleGraph.exists_isBiregularWith_of_lineGraph_regular}
 \uses{def:bipartition,lem:5.4}
 \assume{\begin{pfenum}
  \item $Z$ is a bipartite graph without isolated vertices.
  \item $L(Z)$ is connected and regular.
 \end{pfenum}}
 \prove{$Z$ has a bipartition $(P,Q)$ with $|P|\le|Q|$ that is biregular with some degrees
  $(k_1,k_2)$.}
\end{namedresult}
\begin{proof}
 \pfsketch\ Lemma 5.4 gives a biregular bipartition or regularity, and in a regular bipartite
 graph every bipartition is biregular.
 \step{s5brget}{$Z$ has a biregular bipartition.}
 \begin{proof}
  \step{s5brreg}{\case{$Z$ is $k$-regular.}}
  \begin{proof}\pf\ Any bipartition, which exists by \asm[1]{0}, is biregular with degrees
   $(k,k)$.~\qed\end{proof}
  \step{s5brsemi}{\case{$Z$ has a bipartition that is biregular with degrees $(k_1,k_2)$,
   $k_1\ne k_2$.}}
  \begin{proof}\pf\ This is the claim.~\qed\end{proof}
  \qedstep
  \begin{proof}\pf\ By Lemma~\ref{lem:5.4}, with \asm[1]{0} and \asm[2]{0}, the cases
   \stepref{s5brreg} and \stepref{s5brsemi} are exhaustive.~\qed\end{proof}
 \end{proof}
 \qedstep
 \begin{proof}\pf\ Exchange the two parts of the bipartition of \stepref{s5brget} if
  necessary; this exchanges the two degrees.~\qed\end{proof}
\end{proof}

\begin{cgss}{Lemma}{5.6}\label{lem:5.6}
 \lean{SimpleGraph.lineGraph_cospectral_bipartite_forall_isBiregularWith}
 \uses{def:bipartition,lem:rootgraph,lem:tworoot,lem:biregular,st:5.5,st:6.3,st:6.7}
 \assume{\begin{pfenum}
  \item $X_1$ is a connected bipartite graph with at least one edge, and $L(X_1)$ is
   $h$-regular.
  \item $X_2$ is a graph without isolated vertices.
  \item $L(X_1)$ and $L(X_2)$ are cospectral.
 \end{pfenum}}
 \prove{$X_2$ is connected, and one of the following holds.
  \begin{pfenum}
   \item $X_2$ is bipartite, and for every bipartition $(P_1,Q_1)$ of $X_1$ with
    $|P_1|\le|Q_1|$ that is biregular with degrees $(k_1,k_2)$ there is a bipartition
    $(P_2,Q_2)$ of $X_2$ with $|P_2|=|P_1|$ and $|Q_2|=|Q_1|$ that is biregular with degrees
    $(k_1,k_2)$, such that $\Spec(B_1B_1^T)=\Spec(B_2B_2^T)$ for the biadjacency matrices
    $B_1$, $B_2$ of $X_1$, $X_2$ for these bipartitions.
   \item $X_1$ is not regular; $X_2$ is $\kappa$-regular and not bipartite, for the natural
    number $\kappa$ with $2\kappa=h+2$; and $|V(X_1)|=|V(X_2)|+1$.
  \end{pfenum}}
\end{cgss}
\begin{proof}
 \pfsketch\ \repair\ We need \asm[2]{0} (errata, Implicit hypothesis 5: $X_1=K_{1,3}$ and $X_2$ a
 triangle with an isolated vertex have the same line graph). \new\ We state (i) for every
 biregular bipartition of $X_1$, so that Theorem 5.8 can choose the points of a design as $P_1$;
 the toolkit's Theorem 6.7 orders the parts by size instead, and so this lemma has its own Lean
 name. We will take connectedness and (ii) from the toolkit's Theorem 6.7, which follows the
 paper: in the nonbipartite case $M_2M_2^T$ is nonsingular by Lemma 5.3, so $X_2$ has $p+q-1$
 vertices. For (i) we will follow the first part of the paper's proof, which is
 Lemmas~\ref{lem:rootgraph} and~\ref{lem:tworoot}.
 \step{s556L1}{$X_1$ has no isolated vertices, and $L(X_1)$ is connected and $h$-regular on a
  nonempty vertex set.}
 \begin{proof}\pf\ By \asm[1]{0} and \ST{Lemma 6.3}: a connected graph with an edge has no
  isolated vertices, since each vertex has a walk to an end of the edge that is different from
  it.~\qed\end{proof}
 \step{s556L2}{$L(X_2)$ is connected and $h$-regular.}
 \begin{proof}\pf\ \ST{Corollary 5.5}, with \stepref{s556L1} and \asm[3]{0}.~\qed\end{proof}
 \step{s556st}{$X_2$ has an edge and is connected, and if $X_2$ is not bipartite, then $X_1$ is
  not regular, $X_2$ is regular, and $|V(X_1)|=|V(X_2)|+1$.}
 \begin{proof}\pf\ \ST{Theorem 6.7}, with \asm[1]{0}, \asm[2]{0} and \asm[3]{0}; $L(X_2)$ is
  nonempty by \stepref{s556L2}, so $X_2$ has an edge.~\qed\end{proof}
 \step{s556bip}{\case{$X_2$ is bipartite.}\prove{(i) holds.}}
 \begin{proof}
  \step{s556given}{\assume{$(P_1,Q_1)$ is a bipartition of $X_1$ with $|P_1|\le|Q_1|$ that is
   biregular with degrees $(k_1,k_2)$.}
   \prove{There is a bipartition $(P_2,Q_2)$ of $X_2$ as in (i).}}
  \begin{proof}
   \step{s556pick}{\pick\ a bipartition $(P_2,Q_2)$ of $X_2$ with $|P_2|\le|Q_2|$ that is
    biregular with some degrees $(k_1',k_2')$.}
   \begin{proof}\pf\ Lemma~\ref{lem:biregular} for $X_2$, with \asm{1}, \asm[2]{0}
    and \stepref{s556L2}.~\qed\end{proof}
   \qedstep
   \begin{proof}\pf\ Lemma~\ref{lem:tworoot} with $Z=X_1$, $Z'=X_2$ applies by \asm[1]{0},
    \stepref{s556st}, \asm{2}, \stepref{s556pick} and \asm[3]{0}; it gives the sizes, the
    degrees $k_1'=k_1$, $k_2'=k_2$ and the spectra.~\qed\end{proof}
  \end{proof}
  \qedstep
  \begin{proof}\pf\ \asm{1} and \stepref{s556given} are (i).~\qed\end{proof}
 \end{proof}
 \step{s556nonbip}{\case{$X_2$ is not bipartite.}\prove{(ii) holds.}}
 \begin{proof}
  \step{s556reg}{$X_2$ is $\kappa$-regular for a natural number $\kappa\ge1$ with $2\kappa=h+2$.}
  \begin{proof}\pf\ $X_2$ is regular by \stepref{s556st} and \asm{1}, of some valency $\kappa$;
   $\kappa\ge1$ since $X_2$ has an edge, by \stepref{s556st}. By the degree formula of
   \ST{Lemma 6.3}, with \asm[2]{0}, an edge of $X_2$ has degree $2\kappa-2$ in $L(X_2)$, which is
   $h$ by \stepref{s556L2}.~\qed\end{proof}
  \qedstep
  \begin{proof}\pf\ Steps \stepref{s556st} and \stepref{s556reg}, with \asm{1}, are
   (ii).~\qed\end{proof}
 \end{proof}
 \qedstep
 \begin{proof}\pf\ Step \stepref{s556st}, and the cases \stepref{s556bip} and
  \stepref{s556nonbip}, which are exhaustive.~\qed\end{proof}
\end{proof}

\begin{cgss}{Remark}{5.7}\label{rem:5.7}
 \lean{SimpleGraph.lineGraph_cospectral_bipartite_spectrum_of_not_isBipartite}
 \uses{def:bipartition,lem:5.6,lem:rootgraph,st:6.8}
 \assume{\begin{pfenum}
  \item $X_1$ is a connected graph with at least one edge, and $(P,Q)$ is a bipartition of
   $X_1$, biregular with degrees $(k_1,k_2)$ with $k_1\ne k_2$, with $p=|P|\le q=|Q|$ and
   biadjacency matrix $B$ of rank $\rho$.
  \item $X_2$ is a graph without isolated vertices that is not bipartite, and $L(X_1)$ and
   $L(X_2)$ are cospectral.
 \end{pfenum}}
 \prove{$k_1>k_2\ge1$; $k_1k_2$ is an eigenvalue of $BB^T$ of multiplicity 1; and if
  $\Lambda'$ is the multiset of positive eigenvalues of $BB^T$ with one copy of $k_1k_2$ removed,
  then
  \[
   \Spec(X_2)=\ms{\tfrac12(k_1+k_2)}
    +\sum_{\mu\in\Lambda'}\ms{\delta_\mu,-\delta_\mu}
    +\ms{\bigl(\tfrac12(k_1-k_2)\bigr)^{p-\rho},\bigl(-\tfrac12(k_1-k_2)\bigr)^{q-\rho}},
  \]
  where $\delta_\mu=\sqrt{\mu+\frac14(k_1-k_2)^2}$.}
 We state only the conclusion of the remark, the spectrum of $A(X_2)$. Its intermediate
 assertions, that each $\mu$ gives the eigenvalues $\pm\sqrt\mu$ of $A(X_1)$, $\gamma$ and
 $\gamma'$ of $C$, and $\gamma-2$ and $\gamma'-2$ of $A(L(X_1))$, have no Lean statements of
 their own, and neither has the paper's comment that the multiplicities of $\pm\frac12(k_1-k_2)$
 need not be equal (the statement above gives them, as $p-\rho$ and $q-\rho$). The spectrum of a
 matrix of the form of $C$ is \path{Matrix.spectrumMultiset_fromBlocks_scalar_of_ne}.
\end{cgss}
\begin{proof}
 \pfsketch\ \repair\ The paper says that each eigenvalue $\mu\ne0$ of $B_1B_1^T$ gives two
 eigenvalues $\gamma$, $\gamma'$ of $C$, the roots of $(\gamma-k_1)(\gamma-k_2)=\mu$, and two
 eigenvalues $\delta=\gamma-\frac12(k_1+k_2)$ and $\delta'=\gamma'-\frac12(k_1+k_2)$ of $A(X_2)$
 (errata, Error 13, item 8, for $\delta'$), with $\delta^2=\mu+\frac14(k_1-k_2)^2$. This fails for
 $\mu=k_1k_2$, whose roots are $k_1+k_2$ and $0$, and $0$ is not an eigenvalue of the nonsingular
 $M_2M_2^T$; and the multiplicities of $\pm\frac12(k_1-k_2)$ are $p-\rho$ and $q-\rho$, whose
 difference, not sum, is $q-p$ (errata, Error 14). The toolkit's Corollary 6.8 is the paper's
 computation with these corrections, for $X_2$ connected, regular and not bipartite. We will
 derive those hypotheses from Lemma 5.6; the form that assumes only that $X_2$ has no isolated
 vertices is the one that Theorem 5.8 uses, so this remark has its own Lean name.
 \step{s557deg}{$k_1>k_2\ge1$.}
 \begin{proof}\pf\ By Lemma~\ref{lem:rootgraph}, assertion 1, $k_2\ge1$ and $pk_1=qk_2$, and
  $p\le q$, $k_1\ne k_2$.~\qed\end{proof}
 \step{s557ii}{$X_2$ is connected and regular, and not bipartite.}
 \begin{proof}\pf\ Lemma~\ref{lem:5.6} applies, since $X_1$ is bipartite by \asm[1]{0} and
  $L(X_1)$ is $(k_1+k_2-2)$-regular by Lemma~\ref{lem:rootgraph}, assertion 2; its alternative
  (i) fails since $X_2$ is not bipartite, so (ii) holds.~\qed\end{proof}
 \qedstep
 \begin{proof}\pf\ \ST{Corollary 6.8}, with \asm[1]{0}, \stepref{s557deg}, \stepref{s557ii} and
  \asm[2]{0}; its $\mu_2,\dots,\mu_\rho$ are the members of $\Lambda'$, and its
  $\pm\sqrt{\mu_j+\frac14(k_1-k_2)^2}$ are $\pm\delta_{\mu_j}$.~\qed\end{proof}
\end{proof}

\subsection{Theorem 5.8 and its consequences}

\begin{cgss}{Theorem}{5.8}\label{thm:5.8}
 \lean{Design.flagGraph_cospectral}
 \uses{def:flag,def:bipartition,hc:def:design,lem:design,lem:incidence,lem:dense,lem:isolated,lem:5.1,lem:5.2,lem:5.6,rem:5.7,def:CP,def:represented,thm:4.11,def:3.2,facts:sec5,st:2.6,st:3.2,st:3.4,st:3.5,st:5.1,st:5.5,st:7.7,st:7.10}
 \assume{\begin{pfenum}
  \item $D_1$ is a design with parameters $(v,k,b,r,\lambda)$ on $P$ and $\cB$, with incidence
   matrix $N$, and $b\ge v$.
  \item $X_1=F(D_1)$, and $X_2$ is a graph cospectral with $X_1$.
 \end{pfenum}}
 \prove{One of the following holds.
  \begin{pfenum}
   \item $v+b\le9$, and $X_2$ is represented in $\typ E_8$.
   \item $X_2\cong F(D_2)$ for a design $D_2$ with the parameters $(v,k,b,r,\lambda)$.
   \item There are integers $s\ge2$ and $t\ge1$ with $t(s+1)$ even, $s\ge2t-1$ and
    $(s-1)\mid2t(t-1)$, such that
    \[
     v=\tfrac12s(s+1),\quad k=st,\quad b=\tfrac12(s+1)(s+2),\quad r=(s+2)t,\quad
     \lambda(s-1)=2t(st-1),
    \]
    and $X_2\cong L(Y)$ for a connected $(s+1)t$-regular graph $Y$ on $s(s+2)$ vertices with
    \[
     \Spec(Y)=\ms{\bigl((s+1)t\bigr)^{1},(-t)^{s+1},(\sqrt c)^{\ell},(-\sqrt c)^{\ell}},\qquad
     c=\frac{t(s+1)(s-t)}{s-1},\quad \ell=\tfrac12(s-1)(s+2).
    \]
  \end{pfenum}}
\end{cgss}
\begin{proof}
 \pfsketch\ \repair\ The hypothesis $b\ge v$ is the paper's remark after Theorem 5.8: by
 Lemma~\ref{lem:design} it holds when $\lambda<r$, and otherwise every block contains every point
 and one passes to the dual design, as Corollary 5.9 does. Without it the theorem is false
 (errata, Error 16): the design $(10,10,6,6,6)$ has flag graph $L(K_{6,10})$, cospectral with
 $L(T(6))$, and none of (i), (ii), (iii) holds. Here $v+b=16$; (ii) fails since the only design
 with these parameters has flag graph $L(K_{6,10})$, and $T(6)$ is not bipartite, so Whitney's
 theorem excludes $L(T(6))\cong L(K_{6,10})$; and $v=10$ gives $s=4$ and $t=k/s=\frac52$. The
 details are in Error 16. We will follow the paper. The spectrum makes $X_2$
 regular and connected with $2I+A(X_2)\succeq0$ of rank $v+b-1$, so Theorem 4.11 applies. A
 representation in $\typ E_8$ gives (i), a cocktail party graph gives (ii), and a line graph
 $L(Y)$ gives (ii) or (iii) according to the two cases of Lemma 5.6.
 \step{s558X1}{$X_1$ is connected and $(r+k-2)$-regular on $vr$ vertices, and $2I+A(X_1)$ is
  positive semidefinite of rank $v+b-1$.}
 \begin{proof}\pf\ Lemma~\ref{lem:incidence}, assertions 3 and 4.~\qed\end{proof}
 \step{s558X2}{$X_2$ is connected and $(r+k-2)$-regular on $vr$ vertices, and $2I+A(X_2)$ is
  positive semidefinite of rank $v+b-1$.}
 \begin{proof}
  \step{s558X2reg}{$X_2$ is connected and $(r+k-2)$-regular on $vr$ vertices.}
  \begin{proof}\pf\ \ST{Corollary 5.5} and \ST{Proposition 2.6}, with \stepref{s558X1}; the
   corollary is Lemma~\ref{lem:5.2} and Lemma~\ref{lem:5.1} for both graphs, as in the
   paper.~\qed\end{proof}
  \step{s558X2psd}{$2I+A(X_2)$ is positive semidefinite.}
  \begin{proof}\pf\ By \ST{Corollary 3.4} and \stepref{s558X1}, every eigenvalue of $X_1$ is at
   least $-2$; $X_2$ has the same eigenvalues by \ST{Proposition 2.6}; apply \ST{Corollary 3.4}
   again.~\qed\end{proof}
  \step{s558X2rank}{$\rk(2I+A(X_2))=v+b-1$.}
  \begin{proof}\pf\ \ST{Theorem 3.2} with $f(z)=z+2$ gives the rank from the multiplicities,
   which agree by \ST{Proposition 2.6}; use \stepref{s558X1}.~\qed\end{proof}
  \qedstep
  \begin{proof}\pf\ Steps \stepref{s558X2reg}, \stepref{s558X2psd} and
   \stepref{s558X2rank}.~\qed\end{proof}
 \end{proof}
 \step{s558hrc}{$X_2\cong L(Y)$ for some graph $Y$, or $X_2\cong CP(m)$ for some $m$, or $X_2$ is
  represented in $\typ E_8$.}
 \begin{proof}\pf\ Theorem~\ref{thm:4.11}, with
  \stepref{s558X2}.~\qed\end{proof}
 \step{s558E8}{\case{$X_2$ is represented in $\typ E_8$.}\prove{(i) holds.}}
 \begin{proof}
  \step{s558E8rank}{$\rk(2I+A(X_2))\le8$.}
  \begin{proof}\pf\ $2I+A(X_2)$ is the Gram matrix of a family in $\re^8$
   (Definition~\ref{def:3.2}), so by \ST{Theorem 3.5} its rank is the dimension of a subspace of
   $\re^8$; apply \path{Submodule.finrank_le} and \path{finrank_euclideanSpace_fin}.~\qed\end{proof}
  \qedstep
  \begin{proof}\pf\ $v+b-1\le8$ by \stepref{s558X2} and \stepref{s558E8rank}, and $X_2$ is
   represented in $\typ E_8$ by \asm{1}.~\qed\end{proof}
 \end{proof}
 \step{s558CP}{\case{$X_2\cong CP(m)$.}\prove{(ii) holds, with $D_2=D_1$.}}
 \begin{proof}
  \pfsketch\ \repair\ The paper does not exclude this case of Theorem 4.11 (errata, Gap 15). The
  errata's fill forces $v=k=b=r=2$ and hence case (i), which would also need a representation in
  $\typ E_8$; we reach (ii) instead, since both graphs are then cocktail party graphs.
  \step{s558CPX2}{$X_2$ is $(2m-2)$-regular on $2m$ vertices, and $m\ge1$.}
  \begin{proof}\pf\ $CP(m)$ is, by Definition~\ref{def:CP}, and
   $X_2$ has $vr\ge1$ vertices by \stepref{s558X2}.~\qed\end{proof}
  \step{s558CPX1}{$X_1\cong CP(m)$.}
  \begin{proof}\pf\ By \ST{Corollary 5.5} and \ST{Proposition 2.6}, applied to $X_2$ and its
   cospectral mate $X_1$, $X_1$ is $(2m-2)$-regular on $2m$ vertices by \stepref{s558CPX2}; apply
   Lemma~\ref{lem:dense}, assertion 2.~\qed\end{proof}
  \qedstep
  \begin{proof}\pf\ $X_2\cong CP(m)\cong X_1=F(D_1)$ by \asm{1} and
   \stepref{s558CPX1}.~\qed\end{proof}
 \end{proof}
 \step{s558line}{\case{$X_2\cong L(Y)$.}\prove{(ii) or (iii) holds.}}
 \begin{proof}
  \step{s558Yp}{\pick\ a graph $Y'$ without isolated vertices with $X_2\cong L(Y')$.}
  \begin{proof}\pf\ Lemma~\ref{lem:isolated} and \asm{1}.~\qed\end{proof}
  \step{s558cosp}{$L(I(D_1))$ and $L(Y')$ are cospectral.}
  \begin{proof}\pf\ $L(I(D_1))=X_1$ is cospectral with $X_2$ by \asm[2]{0}, and $X_2$ with
   $L(Y')$ by \stepref{s558Yp} and \ST{Proposition 2.6}; cospectrality is equality of
   characteristic polynomials.~\qed\end{proof}
  \step{s558apply}{Lemma~\ref{lem:5.6} applies to $X_1'=I(D_1)$, with $h=r+k-2$, and
   $X_2'=Y'$; in particular $Y'$ is connected.}
  \begin{proof}\pf\ $I(D_1)$ is connected and bipartite with an edge by
   Lemma~\ref{lem:incidence}, assertions 1 and 2, and $L(I(D_1))=X_1$ is $(r+k-2)$-regular by
   \stepref{s558X1}; with \stepref{s558Yp} and \stepref{s558cosp}.~\qed\end{proof}
  \step{s558bip}{\case{Alternative (i) of Lemma~\ref{lem:5.6} holds.}\prove{(ii) holds.}}
  \begin{proof}
   \step{s558bippick}{\pick\ a bipartition $(P_2,Q_2)$ of $Y'$ with $|P_2|=v$, $|Q_2|=b$,
    biregular with degrees $(r,k)$, whose biadjacency matrix $B_2$ has
    $\Spec(B_2B_2^T)=\Spec(NN^T)$.}
   \begin{proof}\pf\ Apply (i) to the bipartition $(P,\cB)$ of $I(D_1)$, which is biregular with
    degrees $(r,k)$ and has biadjacency matrix $N$ by Lemma~\ref{lem:incidence}, and has
    $|P|=v\le b=|\cB|$ by \asm[1]{0}.~\qed\end{proof}
   \step{s558bipgram}{$B_2B_2^T=(r-\lambda)I+\lambda J$.}
   \begin{proof}\pf\ \ST{Theorem 7.7} for the zero-one $v\times b$ matrix $B_2$, whose row sums
    are $r$ and column sums $k$ by \stepref{s558bippick}; its hypotheses $v>1$, $k>1$, $r>0$,
    $\lambda>0$, $r(k-1)=\lambda(v-1)$ and $\Spec(B_2B_2^T)=\ms{(rk)^{1},(r-\lambda)^{v-1}}$ hold
    by Definition~\ref{hc:def:design}, Lemma~\ref{lem:design} and
    \stepref{s558bippick}.~\qed\end{proof}
   \step{s558bipD2}{The rule $D_2(w)=\{u\in P_2:u\sim w\}$ for $w\in Q_2$ is a design with
    parameters $(v,k,b,r,\lambda)$ on $P_2$ and $Q_2$.}
   \begin{proof}\pf\ By Definition~\ref{hc:def:design}: $|P_2|=v$, $|Q_2|=b$; a block has $k$
    points and a point lies in $r$ blocks by \stepref{s558bippick}; two points lie in
    $(B_2B_2^T)(u,u')=\lambda$ common blocks by \stepref{s558bipgram}.~\qed\end{proof}
   \step{s558bipiso}{$I(D_2)\cong Y'$.}
   \begin{proof}\pf\ The map $P_2\sqcup Q_2\to V(Y')$ that is the inclusion on each part is a
    bijection, and it carries the incidences of $D_2$ onto the edges of $Y'$, since every edge of
    $Y'$ joins $P_2$ to $Q_2$ (Definition~\ref{def:bipartition}).~\qed\end{proof}
   \qedstep
   \begin{proof}\pf\ $F(D_2)=L(I(D_2))\cong L(Y')\cong X_2$ by \stepref{s558bipiso},
    \path{SimpleGraph.Iso.lineGraph} and \stepref{s558Yp}.~\qed\end{proof}
  \end{proof}
  \step{s558nb}{\case{Alternative (ii) of Lemma~\ref{lem:5.6} holds.}\prove{(iii) holds.}}
  \begin{proof}
   \step{s558nbfacts}{$r>k$, and $Y'$ is connected, not bipartite and $\kappa$-regular on
    $v+b-1$ vertices, with $vr$ edges, for the natural number $\kappa$ with $2\kappa=r+k$.}
   \begin{proof}
    \step{s558nbrk}{$r>k$.}
    \begin{proof}\pf\ $I(D_1)$ is not regular by (ii), so $r\ne k$ by
     Lemma~\ref{lem:incidence}; and $r=bk/v\ge k$ by Lemma~\ref{lem:design} and
     \asm[1]{0}.~\qed\end{proof}
    \qedstep
    \begin{proof}\pf\ By (ii), $Y'$ is not bipartite and $\kappa$-regular with
     $2\kappa=h+2=r+k$, and $|V(Y')|=(v+b)-1$; $Y'$ is connected by \stepref{s558apply}; and
     $|E(Y')|=|V(L(Y'))|=|V(X_2)|=vr$ by \stepref{s558Yp} and \stepref{s558X2}.~\qed\end{proof}
   \end{proof}
   \step{s558nbhand}{$(r+k)(v+b-1)=4vr$.}
   \begin{proof}\pf\ \path{SimpleGraph.sum_degrees_eq_twice_card_edges} for $Y'$ gives
    $\kappa(v+b-1)=2vr$ by \stepref{s558nbfacts}; multiply by 2.~\qed\end{proof}
   \step{s558nbv}{$v(r-k)^2=k(r+k)$.}
   \lean{Computations.C11.parameter_equation}
   \begin{proof}\pf\ Multiply \stepref{s558nbhand} by $k$ and substitute $bk=vr$ from
    Lemma~\ref{lem:design}: $(r+k)(vk+vr-k)=4vrk$, which rearranges to the claim (computation
    C11).~\qed\end{proof}
   \step{s558nbt}{$t=\frac12(r-k)$ is a positive integer.}
   \lean{Computations.C11.exists_positive_half_difference}
   \begin{proof}\pf\ $r+k=2\kappa$ is even by \stepref{s558nbfacts}, so $r-k=(r+k)-2k$ is even;
    and $r>k$.~\qed\end{proof}
   \step{s558nbtk}{$2t^2v=k(k+t)$ and $t(b-v)=k+t$; thus $t\mid k$.}
   \lean{Computations.C11.two_mul_sq_mul_points}\lean{Computations.C11.mul_blocks_sub_points}
   \lean{Computations.C11.dvd_block_size}
   \begin{proof}\pf\ Substitute $r=k+2t$ in \stepref{s558nbv} for the first; then
    $b-v=v(r-k)/k=2tv/k=(k+t)/t$ by $bk=vr$ (computation C11). As $t\mid k+t$, also
    $t\mid k$.~\qed\end{proof}
   \pflet{$s=k/t$, a positive integer by \stepref{s558nbtk}.}
   \step{s558nbparams}{$v=\frac12s(s+1)$, $k=st$, $b=\frac12(s+1)(s+2)$ and $r=(s+2)t$.}
   \lean{Computations.C11.parameters}
   \begin{proof}\pf\ Substitute $k=st$ in \stepref{s558nbtk} and in $r=k+2t$ and $bk=vr$
    (computation C11).~\qed\end{proof}
   \step{s558nbs}{$s\ge2$.}
   \lean{Computations.C11.two_le_parameter}
   \begin{proof}\pf\ $\frac12s(s+1)=v>1$ by \stepref{s558nbparams}.~\qed\end{proof}
   \step{s558nblam}{$\lambda(s-1)=2t(st-1)$ and $(s-1)\mid2t(t-1)$.}
   \lean{Computations.C11.lambda_equation}\lean{Computations.C11.lambda_remainder}
   \lean{Computations.C11.parameter_dvd}
   \begin{proof}\pf\ Substituting \stepref{s558nbparams} in $r(k-1)=\lambda(v-1)$ of
    Lemma~\ref{lem:design} gives $(s+2)t(st-1)=\frac12\lambda(s+2)(s-1)$; divide by $s+2$. Then
    $2t(t-1)=(\lambda-2t^2)(s-1)$ (computation C11).~\qed\end{proof}
   \step{s558nbineq}{$s\ge2t-1$.}
   \lean{Computations.C11.two_mul_sub_one_le}
   \begin{proof}\pf\ \repair\ The paper prints $s\le2t-1$ (errata, Error 4). By
    Lemma~\ref{lem:design}, $k\le v$, that is $st\le\frac12s(s+1)$ by \stepref{s558nbparams};
    divide by $s>0$.~\qed\end{proof}
   \step{s558nbY}{$Y'$ is connected and $(s+1)t$-regular on $s(s+2)$ vertices.}
   \begin{proof}\pf\ By \stepref{s558nbfacts} and \stepref{s558nbparams},
    $\kappa=\frac12(r+k)=(s+1)t$ and $v+b-1=(s+1)^2-1=s(s+2)$.~\qed\end{proof}
   \step{s558nbspec}{$\Spec(Y')=\ms{\bigl((s+1)t\bigr)^{1},(-t)^{s+1},(\sqrt c)^{\ell},
    (-\sqrt c)^{\ell}}$.}
   \lean{Computations.C11.spectrum_of_parameters}
   \begin{proof}
    \step{s558nbrem}{$\Spec(Y')=\ms{(s+1)t}+\sum_{\mu\in\Lambda'}\ms{\delta_\mu,-\delta_\mu}
     +\ms{t^{v-\rho},(-t)^{b-\rho}}$, where $\rho=\rk N$, $\delta_\mu=\sqrt{\mu+t^2}$, and
     $\Lambda'$ is the multiset of positive eigenvalues of $NN^T$ with one copy of $rk$ removed.}
    \begin{proof}\pf\ Remark~\ref{rem:5.7} with $X_1=I(D_1)$ and the bipartition $(P,\cB)$,
     biregular with degrees $(r,k)$, $r\ne k$, $v\le b$ by \asm[1]{0}, and biadjacency matrix $N$
     (Lemma~\ref{lem:incidence}), and $X_2=Y'$ by \stepref{s558Yp}, \stepref{s558cosp} and
     \stepref{s558nbfacts}; here $\frac12(r+k)=(s+1)t$ and $\frac12(r-k)=t$.~\qed\end{proof}
    \step{s558nbc}{$r-\lambda+t^2=c$, $v-1=\ell$ and $b-v=s+1$.}
    \lean{Computations.C11.nonprincipal_square}\lean{Computations.C11.spectral_parameters}
    \begin{proof}\pf\ Substitute \stepref{s558nbparams} and \stepref{s558nblam}
     (computation C11).~\qed\end{proof}
    \step{s558nbfish}{\case{$r>\lambda$.}}
    \begin{proof}
     \step{s558nbfishrho}{$\rho=v$ and $\Lambda'=\ms{(r-\lambda)^{v-1}}$.}
     \begin{proof}\pf\ By Lemma~\ref{lem:design}, $\Spec(NN^T)=\ms{(rk)^{1},(r-\lambda)^{v-1}}$
      with both values positive, so $\rk(NN^T)=v$ by \ST{Theorem 3.2}, and $\rho=\rk(NN^T)$ by
      \ST{Theorem 3.5}.~\qed\end{proof}
     \qedstep
     \begin{proof}\pf\ By \stepref{s558nbrem}, \stepref{s558nbfishrho} and \stepref{s558nbc},
      the terms are $(s+1)t$, $(\pm\sqrt c)^{v-1}$ with $v-1=\ell$, $t^{0}$ and
      $(-t)^{b-v}$ with $b-v=s+1$. (Computation C11,
      \path{Computations.C11.spectrum_of_parameters}, is this substitution; it takes the spectrum
      of \stepref{s558nbrem} in this form as a hypothesis, so the rank cases and
      Remark~\ref{rem:5.7} stay with this proof.)~\qed\end{proof}
    \end{proof}
    \step{s558nbtriv}{\case{$r=\lambda$.}}
    \begin{proof}
     \step{s558nbtrivrho}{$\rho=1$ and $\Lambda'$ is empty.}
     \begin{proof}\pf\ By Lemma~\ref{lem:design}, $\Spec(NN^T)=\ms{(rk)^{1},0^{v-1}}$ with
      $rk>0$, so $\rk(NN^T)=1$ by \ST{Theorem 3.2}, and $\rho=\rk(NN^T)$ by
      \ST{Theorem 3.5}.~\qed\end{proof}
     \step{s558nbtrivc}{$c=t^2$, so $\sqrt c=t$.}
     \begin{proof}\pf\ $c=r-\lambda+t^2$ by \stepref{s558nbc}, and $r=\lambda$.~\qed\end{proof}
     \qedstep
     \begin{proof}\pf\ By \stepref{s558nbrem} and \stepref{s558nbtrivrho},
      $\Spec(Y')=\ms{(s+1)t,t^{v-1},(-t)^{b-1}}$, and $v-1=\ell$, $b-1=(s+1)+\ell$ by
      \stepref{s558nbc}; with \stepref{s558nbtrivc} this is the claim.~\qed\end{proof}
    \end{proof}
    \qedstep
    \begin{proof}\pf\ $r\ge\lambda$ by Lemma~\ref{lem:design}, so the cases
     \stepref{s558nbfish} and \stepref{s558nbtriv} are exhaustive.~\qed\end{proof}
   \end{proof}
   \step{s558nbeven}{$t(s+1)$ is even.}
   \lean{Computations.C11.even_cube_diagonal}\lean{Computations.C11.even_parameter_cube}
   \lean{Computations.C11.even_mul_of_even_cube}
   \begin{proof}
    \pflet{$A=A(Y')$, $f(z)=(z+t)(z^2-c)$ and $e=t^2(s+1)(st-1)/(s-1)$.}
    \step{s558nbhp}{$f(A)=eJ$.}
    \lean{Computations.C11.matrix_identity}\lean{Computations.C11.hoffman_value}
    \begin{proof}\pf\ $f$ vanishes at $-t$ and $\pm\sqrt c$, which by \stepref{s558nbspec} are
     the eigenvalues other than $(s+1)t$; \ST{Theorem 7.10} with \stepref{s558nbY} gives
     $f(A)=\frac{f((s+1)t)}{s(s+2)}J$, and $f((s+1)t)=s(s+2)e$ (computation C11).~\qed\end{proof}
    \step{s558nbdiag}{$A^3(u,u)=t^3(s+1)$ for every vertex $u$.}
    \lean{Computations.C11.cube_diagonal_eq}\lean{Computations.C11.triangle_diagonal}
    \begin{proof}\pf\ The $(u,u)$ entry of $f(A)=A^3+tA^2-cA-ctI$ is
     $A^3(u,u)+t^2(s+1)-ct$, since $A(u,u)=0$ and $A^2(u,u)=\deg(u)=(s+1)t$ by
     \path{SimpleGraph.adjMatrix_mul_self_apply_self};
     by \stepref{s558nbhp} it equals $e$, and $e+ct-t^2(s+1)=t^3(s+1)$ (computation
     C11).~\qed\end{proof}
    \qedstep
    \begin{proof}\pf\ $Y'$ has a vertex, and $A^3(u,u)$ is twice a number of triangles by
     \ST{Theorem 5.1}; so $t^3(s+1)$ is even by \stepref{s558nbdiag}. The difference
     $t^3(s+1)-t(s+1)=(t-1)t(t+1)(s+1)$ is even, as a multiple of two consecutive
     integers. (In computation C11, \path{Computations.C11.even_parameter_cube} takes the
     integer form of \stepref{s558nbdiag} as a hypothesis; the passage from the real identity of
     \stepref{s558nbhp} to it stays with this proof.)~\qed\end{proof}
   \end{proof}
   \qedstep
   \begin{proof}\pf\ Take $Y=Y'$: $X_2\cong L(Y')$ by \stepref{s558Yp}; the conditions are
    \stepref{s558nbparams} to \stepref{s558nbeven}, with $t\ge1$ by \stepref{s558nbt}.~\qed\end{proof}
  \end{proof}
  \qedstep
  \begin{proof}\pf\ By \stepref{s558apply}, the cases \stepref{s558bip} and \stepref{s558nb} are
   the two alternatives of Lemma~\ref{lem:5.6}.~\qed\end{proof}
 \end{proof}
 \qedstep
 \begin{proof}\pf\ By \stepref{s558hrc}, the cases \stepref{s558E8}, \stepref{s558CP} and
  \stepref{s558line} are exhaustive.~\qed\end{proof}
\end{proof}

\medskip\noindent\textbf{Corollaries 5.9 to 5.11 and Remark 5.12.}

The next two lemmas settle the small cases that the corollaries leave to the reader.

\begin{namedresult}{Lemma}{(a single block)}\label{lem:oneblock}
 \lean{Design.flagGraphCharacterized_of_b_eq_one}
 \lean{SimpleGraph.Cospectral.nonempty_iso_of_isRegularOfDegree_card_sub_one}
 \uses{def:flag,lem:incidence,lem:dense,st:2.6,st:5.5}
 \prove{\begin{pfenum}
  \item If $X$ is an $(n-1)$-regular graph on $n\ge1$ vertices, then every graph cospectral
   with $X$ is isomorphic to $X$.
  \item If $D$ is a design with parameters $(v,v,1,1,1)$, then $F(D)$ is characterized by its
   spectrum.
 \end{pfenum}}
\end{namedresult}
\begin{proof}
 \pfsketch\ \repair\ The paper's Corollary 5.11 omits the designs with a single block (errata,
 Gap 16). A cospectral mate is regular of the same valency on the same number of vertices, so
 both graphs are complete.
 \step{s5obmate}{\assume{$X$ is an $(n-1)$-regular graph on $n\ge1$ vertices, and $X'$ is
  cospectral with $X$.}\prove{$X'\cong X$.}}
 \begin{proof}\pf\ $X'$ is $(n-1)$-regular on $n$ vertices by \ST{Corollary 5.5} and
  \ST{Proposition 2.6}; both graphs are isomorphic to $K_n$ by Lemma~\ref{lem:dense},
  assertion 1.~\qed\end{proof}
 \step{s5obflag}{\assume{$D$ is a design with parameters $(v,v,1,1,1)$.}
  \prove{$F(D)$ is characterized by its spectrum.}}
 \begin{proof}
  \step{s5obflagreg}{$F(D)$ is $(v-1)$-regular on $v$ vertices.}
  \begin{proof}\pf\ Lemma~\ref{lem:incidence}, assertion 3, with $r+k-2=v-1$ and
   $vr=v$.~\qed\end{proof}
  \qedstep
  \begin{proof}\pf\ Apply \stepref{s5obmate} to $F(D)$, by \stepref{s5obflagreg}, and take
   $D_2=D$ in Definition~\ref{def:flag}.~\qed\end{proof}
 \end{proof}
 \qedstep
 \begin{proof}\pf\ Steps \stepref{s5obmate} and \stepref{s5obflag}.~\qed\end{proof}
\end{proof}

\begin{namedresult}{Lemma}{(the hexagon)}\label{lem:hexagon}
 \lean{Design.flagGraphCharacterized_of_parameters_3_2_1}
 \uses{def:flag,lem:incidence,lem:cycle,st:2.6,st:5.5}
 \assume{$D$ is a design with parameters $(3,2,3,2,1)$.}
 \prove{$F(D)$ is characterized by its spectrum.}
\end{namedresult}
\begin{proof}
 \pfsketch\ \repair\ The paper leaves this to the reader in Corollaries 5.10 and 5.11 (errata,
 Gap 16): a cospectral mate of the hexagon is connected and $2$-regular on six vertices.
 \step{s5hxflag}{$F(D)\cong C_6$.}
 \begin{proof}\pf\ $F(D)$ is connected and $2$-regular on $6$ vertices by
  Lemma~\ref{lem:incidence}, assertion 3; apply Lemma~\ref{lem:cycle}.~\qed\end{proof}
 \step{s5hxmate}{\assume{$X_2$ is cospectral with $F(D)$.}\prove{$X_2\cong F(D)$.}}
 \begin{proof}\pf\ $X_2$ is connected and $2$-regular on $6$ vertices by \ST{Corollary 5.5} and
  \ST{Proposition 2.6}, so $X_2\cong C_6$ by Lemma~\ref{lem:cycle}; use
  \stepref{s5hxflag}.~\qed\end{proof}
 \qedstep
 \begin{proof}\pf\ Step \stepref{s5hxmate} with $D_2=D$ in Definition~\ref{def:flag}.~\qed\end{proof}
\end{proof}

\begin{namedresult}{Lemma}{(designs with full blocks)}\label{lem:fullblock}
 \lean{Design.trivial}\lean{Design.nonempty_iso_flagGraph_lineGraph_completeBipartite}
 \uses{def:flag,hc:def:design,lem:design,facts:sec5}
 \assume{$m\ge1$ and $n\ge2$ are integers.}
 \prove{\begin{pfenum}
  \item The \emph{trivial design} $T$ on the points $\{1,\dots,n\}$ with $m$ block labels, each
   block being the whole point set, is a design with parameters $(n,n,m,m,m)$.
  \item If $D$ is a design with parameters $(n,n,m,m,m)$, then $F(D)\cong L(K_{m,n})$.
 \end{pfenum}}
\end{namedresult}
\begin{proof}
 \pfsketch\ Every point lies in every block, so the incidence graph is complete bipartite.
 \step{s5fbT}{Assertion 1 holds.}
 \begin{proof}\pf\ By Definition~\ref{hc:def:design}: there are $n>1$ points and $m>0$ block
  labels, every block has $n$ points, and every point and every pair of points lies in all $m$
  blocks.~\qed\end{proof}
 \step{s5fbiso}{\assume{$D$ is a design with parameters $(n,n,m,m,m)$ on $P$ and
  $\cB$.}\prove{$F(D)\cong L(K_{m,n})$.}}
 \begin{proof}
  \step{s5fbfull}{Every block of $D$ contains every point.}
  \begin{proof}\pf\ Lemma~\ref{lem:design}, assertion 2, since $k=v=n$.~\qed\end{proof}
  \step{s5fbI}{$I(D)\cong K_{m,n}$.}
  \begin{proof}\pf\ By \stepref{s5fbfull}, every point is adjacent to every block in $I(D)$,
   so $I(D)$ is the complete bipartite graph on $P\sqcup\cB$ (Definition~\ref{def:flag}), with
   $|P|=n$ and $|\cB|=m$. We obtain an isomorphism with $K_{m,n}$ from \path{Equiv.sumComm},
   \path{Fintype.equivFin} and \path{SimpleGraph.completeBipartiteGraphCongr}.~\qed\end{proof}
  \qedstep
  \begin{proof}\pf\ $F(D)=L(I(D))$, and apply \path{SimpleGraph.Iso.lineGraph} to
   \stepref{s5fbI}.~\qed\end{proof}
 \end{proof}
 \qedstep
 \begin{proof}\pf\ Steps \stepref{s5fbT} and \stepref{s5fbiso}.~\qed\end{proof}
\end{proof}

\begin{cgss}{Corollary}{5.9}\label{cor:5.9}
 \lean{SimpleGraph.lineGraph_completeBipartite_characterized}
 \uses{thm:5.8,def:flag,hc:def:design,lem:design,lem:oneblock,lem:fullblock,facts:sec5,st:2.6,st:6.3,st:7.10}
 \assume{\begin{pfenum}
  \item $m,n\ge1$ are integers with $m+n>9$.
  \item There is no integer $t\ge1$ with $\{m,n\}=\{2t^2+t,2t^2-t\}$ for which a symmetric
   Hadamard matrix of order $4t^2$ with constant diagonal exists.
  \item $X_2$ is a graph cospectral with $L(K_{m,n})$.
 \end{pfenum}}
 \prove{$X_2\cong L(K_{m,n})$.}
\end{cgss}
\begin{proof}
 \pfsketch\ \repair\ The two readings of ``characterized'' agree here (errata, Implicit
 hypothesis 6), since a design with $k=v$ is unique. \repair\ We write the exception as a
 condition on $\{m,n\}$ (errata, Error 17): as printed, with $m=6$, $n=10$ the exception does not
 apply, but $L(K_{6,10})$ is cospectral with $L(T(6))$ (errata, Error 14, evidence). We will follow the paper: $L(K_{m,n})$
 is the flag graph of the trivial design, and in case (iii) of Theorem 5.8 bordering a
 $(\pm1)$ form of $A(Y)$ gives the Hadamard matrix.
 \step{s559ordered}{\assume{$m\ge n$, and \asm[1]{0} to \asm[3]{0}.}\prove{$X_2\cong
  L(K_{m,n})$.}}
 \begin{proof}
  \step{s559one}{\case{$n=1$.}}
  \begin{proof}\pf\ $L(K_{m,1})$ is $(m-1)$-regular on $m$ vertices by the degree formula of
   \ST{Lemma 6.3}; apply Lemma~\ref{lem:oneblock}.~\qed\end{proof}
  \step{s559two}{\case{$n\ge2$.}}
  \begin{proof}
   \pflet{$T$ be the trivial design of Lemma~\ref{lem:fullblock} on $n$ points with $m$ block
    labels.}
   \step{s559T}{$T$ is a design with parameters $(n,n,m,m,m)$, $b=m\ge n=v$, and
    $F(T)\cong L(K_{m,n})$.}
   \begin{proof}\pf\ Lemma~\ref{lem:fullblock}, both assertions, with $n\ge2$ by \asm{2} and
    $m\ge n$ by \asm{1}.~\qed\end{proof}
   \step{s559cases}{One of (i), (ii) and (iii) of Theorem~\ref{thm:5.8} holds for $D_1=T$ and
    $X_2$.}
   \begin{proof}\pf\ $X_2$ is cospectral with $F(T)$ by \asm[3]{0}, \stepref{s559T} and
    \ST{Proposition 2.6}; apply Theorem~\ref{thm:5.8}.~\qed\end{proof}
   \step{s559i}{\case{(i) holds.}}
   \begin{proof}\pf\ Then $n+m=v+b\le9$, against \asm[1]{0}.~\qed\end{proof}
   \step{s559ii}{\case{(ii) holds.}}
   \begin{proof}\pf\ Then $X_2\cong F(D_2)$ for a design $D_2$ with parameters $(n,n,m,m,m)$,
    and $F(D_2)\cong L(K_{m,n})$ by Lemma~\ref{lem:fullblock}, assertion 2.~\qed\end{proof}
   \step{s559iii}{\case{(iii) holds.}\prove{\textsc{False}.}}
   \begin{proof}
    \pfsketch\ We will show that the exception of \asm[2]{0} occurs.
    \step{s559st}{$s=2t-1$, $n=2t^2-t$, $m=2t^2+t$, $t\ge2$ and $c=t^2$.}
    \lean{Computations.C12.full_block_parameters}
    \begin{proof}\pf\ From $v=k$ in (iii), $\frac12s(s+1)=st$, so $s+1=2t$; substitute in
     $v$, $b$ and $c$, and $2t-1=s\ge2$ (computation C12). (Computation C12 proves the whole
     step, $t\ge2$ and $c=t^2$ included, for the integers $s\ge2$ and $t\ge1$ of (iii), with the
     fractions evaluated over $\rats$.)~\qed\end{proof}
    \step{s559Y}{$A=A(Y)$ satisfies $A^2=t^2(I+J)$ and $AJ=2t^2J$, on $4t^2-1$ vertices.}
    \begin{proof}\pf\ By (iii) and \stepref{s559st}, $Y$ is connected and $2t^2$-regular on
     $s(s+2)=4t^2-1$ vertices with eigenvalues $2t^2$ and $\pm t$. \ST{Theorem 7.10} with
     $f(z)=z^2-t^2$ gives $A^2-t^2I=\frac{4t^4-t^2}{4t^2-1}J=t^2J$.~\qed\end{proof}
    \pflet{$S=2A-J$ and $H=\pmat{-1&-\Zj^T\\-\Zj&S}$, of order $4t^2$.}
    \lean{Computations.C12.core}\lean{Computations.C12.border}
    \step{s559S}{$S$ is symmetric with entries $\pm1$ and diagonal $-1$, $S\Zj=\Zj$ and
     $S^2=4t^2I-J$.}
    \lean{Computations.C12.core_row_sum}\lean{Computations.C12.core_mul}
    \begin{proof}\pf\ $A$ is a symmetric zero-one matrix with zero diagonal; by \stepref{s559Y},
     $S\Zj=(4t^2-(4t^2-1))\Zj$ and $S^2=4A^2-2AJ-2JA+J^2=4t^2I-J$, since $J^2=(4t^2-1)J$
     (computation C12).~\qed\end{proof}
    \step{s559H}{$H$ is a symmetric Hadamard matrix with constant diagonal $-1$.}
    \lean{Computations.C12.border_isHadamard}\lean{Computations.C12.hadamard_border}
    \begin{proof}\pf\ By \stepref{s559S}, $H$ is symmetric with entries $\pm1$ and diagonal
     $-1$, and block multiplication gives $H^2=4t^2I$: the corner is $1+(4t^2-1)$, the border
     is $\Zj-S\Zj=0$, and the rest is $J+S^2$. This is \path{Matrix.IsHadamard} over
     $\re$.~\qed\end{proof}
    \qedstep
    \begin{proof}\pf\ \stepref{s559st} and \stepref{s559H} contradict \asm[2]{0}.~\qed\end{proof}
   \end{proof}
   \qedstep
   \begin{proof}\pf\ By \stepref{s559cases}, the cases \stepref{s559i}, \stepref{s559ii} and
    \stepref{s559iii} are exhaustive, and \stepref{s559i} and \stepref{s559iii} are
    impossible.~\qed\end{proof}
  \end{proof}
  \qedstep
  \begin{proof}\pf\ The cases \stepref{s559one} and \stepref{s559two} are exhaustive, since
   $n\ge1$.~\qed\end{proof}
 \end{proof}
 \step{s559ge}{\case{$m\ge n$.}}
 \begin{proof}\pf\ Step \stepref{s559ordered}.~\qed\end{proof}
 \step{s559lt}{\case{$m<n$.}}
 \begin{proof}
  \step{s559swap}{$L(K_{m,n})\cong L(K_{n,m})$.}
  \begin{proof}\pf\ \path{Equiv.sumComm} is an isomorphism $K_{m,n}\cong K_{n,m}$ by the
   definition of the complete bipartite graph; apply \path{SimpleGraph.Iso.lineGraph}.~\qed\end{proof}
  \qedstep
  \begin{proof}\pf\ \asm[1]{0} and \asm[2]{0} are symmetric in $m$ and $n$, and $X_2$ is
   cospectral with $L(K_{n,m})$ by \stepref{s559swap} and \ST{Proposition 2.6}; apply
   \stepref{s559ordered} to $(n,m)$, and then \stepref{s559swap}.~\qed\end{proof}
 \end{proof}
 \qedstep
 \begin{proof}\pf\ The cases \stepref{s559ge} and \stepref{s559lt} are exhaustive.~\qed\end{proof}
\end{proof}

\begin{cgss}{Corollary}{5.10}\label{cor:5.10}
 \lean{Design.IsSymmetric.flagGraphCharacterized}
 \uses{thm:5.8,def:flag,hc:def:design,lem:design,lem:hexagon}
 \assume{\begin{pfenum}
  \item $D_1$ is a design with parameters $(v,k,v,r,\lambda)$, that is, a symmetric design.
  \item $k<v$ and $(v,k,\lambda)\ne(4,3,2)$.
 \end{pfenum}}
 \prove{$F(D_1)$ is characterized by its spectrum.}
\end{cgss}
\begin{proof}
 \pfsketch\ We will follow the paper: case (iii) of Theorem 5.8 cannot occur, and case (i)
 leaves only $(3,2,1)$, the hexagon, since $(4,3,2)$ is excluded.
 \step{s5510rk}{$r=k$ and $b\ge v$.}
 \begin{proof}\pf\ $vk=vr$ by Lemma~\ref{lem:design} with $b=v$.~\qed\end{proof}
 \step{s5510mate}{\assume{$X_2$ is cospectral with $F(D_1)$.}
  \prove{$X_2\cong F(D_2)$ for a design $D_2$ with the parameters of $D_1$.}}
 \begin{proof}
  \step{s5510i}{\case{(i) of Theorem~\ref{thm:5.8} holds.}}
  \begin{proof}
   \step{s5510small}{$(v,k,\lambda)=(3,2,1)$.}
   \begin{proof}
    \step{s5510v}{$v\in\{3,4\}$.}
    \begin{proof}\pf\ $2v=v+b\le9$, and $v>k>1$.~\qed\end{proof}
    \step{s5510three}{\case{$v=3$.}}
    \begin{proof}\pf\ $k=2$, and $2\lambda=\lambda(v-1)=r(k-1)=2$ by Lemma~\ref{lem:design}
     and \stepref{s5510rk}, so $\lambda=1$.~\qed\end{proof}
    \step{s5510four}{\case{$v=4$.}\prove{\textsc{False}.}}
    \begin{proof}
     \step{s5510fourlam}{$3\lambda=k(k-1)$ and $k\in\{2,3\}$.}
     \begin{proof}\pf\ $\lambda(v-1)=r(k-1)$ by Lemma~\ref{lem:design}, with $r=k$ by
      \stepref{s5510rk}; and $1<k<v=4$.~\qed\end{proof}
     \step{s5510fourtwo}{\case{$k=2$.}}
     \begin{proof}\pf\ Then $3\lambda=2$ by \stepref{s5510fourlam}, which has no integer
      solution.~\qed\end{proof}
     \step{s5510fourthree}{\case{$k=3$.}}
     \begin{proof}\pf\ Then $\lambda=2$ by \stepref{s5510fourlam}, and $(v,k,\lambda)=(4,3,2)$
      is excluded by \asm[2]{0}.~\qed\end{proof}
     \qedstep
     \begin{proof}\pf\ The cases \stepref{s5510fourtwo} and \stepref{s5510fourthree} are
      exhaustive by \stepref{s5510fourlam}.~\qed\end{proof}
    \end{proof}
    \qedstep
    \begin{proof}\pf\ The cases \stepref{s5510three} and \stepref{s5510four} are exhaustive by
     \stepref{s5510v}.~\qed\end{proof}
   \end{proof}
   \qedstep
   \begin{proof}\pf\ Lemma~\ref{lem:hexagon} with \stepref{s5510small} and $b=v=3$,
    $r=k=2$.~\qed\end{proof}
  \end{proof}
  \step{s5510ii}{\case{(ii) of Theorem~\ref{thm:5.8} holds.}}
  \begin{proof}\pf\ This is the claim.~\qed\end{proof}
  \step{s5510iii}{\case{(iii) of Theorem~\ref{thm:5.8} holds.}\prove{\textsc{False}.}}
  \begin{proof}\pf\ $r=(s+2)t>st=k$ since $t\ge1$, against \stepref{s5510rk}.~\qed\end{proof}
  \qedstep
  \begin{proof}\pf\ Theorem~\ref{thm:5.8} applies by \asm[1]{0}, \stepref{s5510rk} and
   \asm{1}, so the cases are exhaustive.~\qed\end{proof}
 \end{proof}
 \qedstep
 \begin{proof}\pf\ Step \stepref{s5510mate} is Definition~\ref{def:flag}.~\qed\end{proof}
\end{proof}

\begin{remark}
 \lean{Computations.C13.graph1}\lean{Computations.C13.graph2}\lean{Computations.C13.spectrum_graph1}
 \lean{Computations.C13.spectrum_graph2}
 \lean{Computations.C13.cospectral}\lean{Computations.C13.not_nonempty_iso}
 \lean{Computations.C13.design}\lean{Computations.C13.nonempty_iso_flagGraph}
 \lean{Computations.C13.representation_graph1}\lean{Computations.C13.representation_graph2}
 \lean{Computations.C13.isRepresentedIn_E₇_graph1}\lean{Computations.C13.isRepresentedIn_E₇_graph2}
 \lean{Computations.C13.cube}\lean{Computations.C13.cubeLineGraphIso}
 \lean{Computations.C13.nonempty_iso_cube_lineGraph}
 \lean{Computations.C13.not_flagGraphCharacterized}
 The exception is genuine. The paper's two sets of 12 vectors in $\typ E_7$ give two graphs with
 spectrum $4^1\,2^3\,0^3\,(-2)^5$: the flag graph of the $(4,3,2)$ design, which is
 $L(Q_3)$, and a second graph, which computation C13 shows is not isomorphic to the first
 (\path{cgss/prerequisites.md}, HR1). A design with parameters $(4,3,4,3,2)$ is unique: each of
 its four blocks omits one point, each point is omitted by $b-r=1$ block, and so the blocks are
 the four $3$-subsets of the points. Every flag graph of such a design is therefore isomorphic to
 the first graph, and the second graph shows, in the sense of Definition~\ref{def:flag}, that the
 flag graph is not characterized by its spectrum. The uniqueness of the exception (HR2) is a side
 claim, and we do not formalize it here. Computation C13 checks the spectra, that the two graphs are
 represented by the paper's vectors in $\typ E_7$, and that they are not isomorphic (the first
 has no induced claw and the second has one); it identifies the first graph with the line graph
 of the incidence graph of the design on the four $3$-subsets, built in
 \path{LineGraphs/Computations/C13Design.lean} rather than through Definition~\ref{def:flag}.
 It also identifies the first graph with $L(Q_3)$, by an explicit isomorphism, where $Q_3$ has
 the eight binary triples as vertices, two of them adjacent when they differ in exactly one
 coordinate; we do not use this identification. The non-characterization itself is
 \path{Computations.C13.not_flagGraphCharacterized}: the second graph is cospectral with the
 flag graph of the design and has an induced claw, while the flag graph of every design, being a
 line graph, has none.
\end{remark}

\begin{cgss}{Corollary}{5.11}\label{cor:5.11}
 \lean{Design.flagGraphCharacterized_of_lam_eq_one}
 \uses{thm:5.8,def:flag,hc:def:design,lem:design,lem:oneblock,lem:hexagon}
 \assume{$D_1$ is a design with parameters $(v,k,b,r,1)$.}
 \prove{$F(D_1)$ is characterized by its spectrum.}
\end{cgss}
\begin{proof}
 \pfsketch\ \repair\ The paper says that case (i) of Theorem 5.8 occurs only for $(3,2,1)$; it
 also occurs for a single block containing every point, $(v,v,1,1,1)$ with $2\le v\le8$, whose
 flag graph is complete (errata, Gap 16). We will treat that case first, and then follow the
 paper.
 \step{s5511full}{\case{$k=v$.}}
 \begin{proof}\pf\ Then $r=\lambda=1$ by Lemma~\ref{lem:design}, and $b=vr/k=1$, so $D_1$ has
  parameters $(v,v,1,1,1)$; apply Lemma~\ref{lem:oneblock}.~\qed\end{proof}
 \step{s5511proper}{\case{$k<v$.}}
 \begin{proof}
  \step{s5511fisher}{$r>1$ and $b\ge v$.}
  \begin{proof}\pf\ Lemma~\ref{lem:design}: $r\ge\lambda=1$ with equality only if $k=v$, and
   then $b\ge v$ since $r>\lambda$.~\qed\end{proof}
  \step{s5511mate}{\assume{$X_2$ is cospectral with $F(D_1)$.}
   \prove{$X_2\cong F(D_2)$ for a design $D_2$ with the parameters of $D_1$.}}
  \begin{proof}
   \step{s5511i}{\case{(i) of Theorem~\ref{thm:5.8} holds.}}
   \begin{proof}
    \step{s5511small}{$(v,k,b,r)=(3,2,3,2)$.}
    \begin{proof}
     \step{s5511v}{$v\in\{3,4\}$.}
     \begin{proof}\pf\ $2v\le v+b\le9$ by \stepref{s5511fisher}, and $v>k>1$.~\qed\end{proof}
     \step{s5511three}{\case{$v=3$.}}
     \begin{proof}\pf\ $k=2$, $r=\lambda(v-1)/(k-1)=2$ and $b=vr/k=3$ by
      Lemma~\ref{lem:design}.~\qed\end{proof}
     \step{s5511four}{\case{$v=4$.}\prove{\textsc{False}.}}
     \begin{proof}
      \step{s5511fourr}{$r(k-1)=3$ and $k\in\{2,3\}$.}
      \begin{proof}\pf\ $r(k-1)=\lambda(v-1)$ by Lemma~\ref{lem:design}, with $\lambda=1$ and
       $1<k<v=4$.~\qed\end{proof}
      \step{s5511fourtwo}{\case{$k=2$.}}
      \begin{proof}\pf\ Then $r=3$ by \stepref{s5511fourr}, and $b=vr/k=6$ by
       Lemma~\ref{lem:design}, so $v+b=10>9$, against (i).~\qed\end{proof}
      \step{s5511fourthree}{\case{$k=3$.}}
      \begin{proof}\pf\ Then $2r=3$ by \stepref{s5511fourr}, which has no integer
       solution.~\qed\end{proof}
      \qedstep
      \begin{proof}\pf\ The cases \stepref{s5511fourtwo} and \stepref{s5511fourthree} are
       exhaustive by \stepref{s5511fourr}.~\qed\end{proof}
     \end{proof}
     \qedstep
     \begin{proof}\pf\ The cases \stepref{s5511three} and \stepref{s5511four} are exhaustive by
      \stepref{s5511v}.~\qed\end{proof}
    \end{proof}
    \qedstep
    \begin{proof}\pf\ Lemma~\ref{lem:hexagon} with \stepref{s5511small}.~\qed\end{proof}
   \end{proof}
   \step{s5511ii}{\case{(ii) of Theorem~\ref{thm:5.8} holds.}}
   \begin{proof}\pf\ This is the claim.~\qed\end{proof}
   \step{s5511iii}{\case{(iii) of Theorem~\ref{thm:5.8} holds.}\prove{\textsc{False}.}}
   \begin{proof}\pf\ Then $\lambda(s-1)=2t(st-1)\ge2(s-1)$, since $t\ge1$ and $st-1\ge s-1$, so
    $\lambda\ge2$, as $s\ge2$; this contradicts $\lambda=1$.~\qed\end{proof}
   \qedstep
   \begin{proof}\pf\ Theorem~\ref{thm:5.8} applies by \stepref{s5511fisher} and \asm{1}, so the
    cases are exhaustive.~\qed\end{proof}
  \end{proof}
  \qedstep
  \begin{proof}\pf\ Step \stepref{s5511mate} is Definition~\ref{def:flag}.~\qed\end{proof}
 \end{proof}
 \qedstep
 \begin{proof}\pf\ $k\le v$ by Lemma~\ref{lem:design}, so the cases \stepref{s5511full} and
  \stepref{s5511proper} are exhaustive.~\qed\end{proof}
\end{proof}

\begin{cgss}{Remark}{5.12}\label{rem:5.12}
 \lean{Design.flagGraphCharacterized_of_lam_lt_nineteen}
 \uses{thm:5.8,def:flag,hc:def:design,lem:design,hc:cor:complement}
 \assume{$D_1$ is a design with parameters $(v,k,b,r,\lambda)$, with $k<v$ and $2<\lambda<19$.}
 \prove{$F(D_1)$ is characterized by its spectrum.}
\end{cgss}
\begin{proof}
 \pfsketch\ We will follow the paper. In case (iii) of Theorem 5.8,
 $\lambda=2t^2+2t(t-1)/(s-1)$, so $t=1$ gives $\lambda=2$, $t\ge3$ gives $\lambda\ge19$, and
 $t=2$ gives $s\in\{3,5\}$; $s=3$ is a design with $k=v$, and $s=5$ gives the parameters
 $(15,10,21,14,9)$, whose complement $(15,5,21,7,2)$ does not exist. The paper cites Hall's book
 for this; we use the Hall--Connor section. \new\ The paper does not mention case (i); it cannot
 occur for $\lambda\ge3$.
 \step{s5512fisher}{$r>\lambda$ and $b\ge v$.}
 \begin{proof}\pf\ Lemma~\ref{lem:design}, since $k<v$.~\qed\end{proof}
 \step{s5512mate}{\assume{$X_2$ is cospectral with $F(D_1)$.}
  \prove{$X_2\cong F(D_2)$ for a design $D_2$ with the parameters of $D_1$.}}
 \begin{proof}
  \step{s5512i}{\case{(i) of Theorem~\ref{thm:5.8} holds.}\prove{\textsc{False}.}}
  \begin{proof}
   \step{s5512b}{$v\in\{3,4\}$ and $b=\lambda v(v-1)/(k(k-1))$.}
   \begin{proof}\pf\ $2v\le v+b\le9$ by \stepref{s5512fisher}, and $v>k>1$; the formula is
    Lemma~\ref{lem:design}.~\qed\end{proof}
   \step{s5512three}{\case{$v=3$.}}
   \begin{proof}\pf\ $k=2$, so $b=3\lambda\ge9$ and $v+b\ge12$.~\qed\end{proof}
   \step{s5512four}{\case{$v=4$.}}
   \begin{proof}\pf\ $k\in\{2,3\}$, so $b\ge12\lambda/6=2\lambda\ge6$ and
    $v+b\ge10$.~\qed\end{proof}
   \qedstep
   \begin{proof}\pf\ The cases \stepref{s5512three} and \stepref{s5512four} are exhaustive by
    \stepref{s5512b}, and both contradict $v+b\le9$.~\qed\end{proof}
  \end{proof}
  \step{s5512ii}{\case{(ii) of Theorem~\ref{thm:5.8} holds.}}
  \begin{proof}\pf\ This is the claim.~\qed\end{proof}
  \step{s5512iii}{\case{(iii) of Theorem~\ref{thm:5.8} holds.}\prove{\textsc{False}.}}
  \begin{proof}
   \step{s5512lam}{$\lambda(s-1)=2t^2(s-1)+2t(t-1)$.}
   \lean{Computations.C14.lambda_remainder}
   \begin{proof}\pf\ $2t(st-1)=2t^2(s-1)+2t(t-1)$, and (iii) (computation C14).~\qed\end{proof}
   \step{s5512t1}{\case{$t=1$.}}
   \lean{Computations.C14.lambda_eq_two}
   \begin{proof}\pf\ Then $\lambda=2$ by \stepref{s5512lam}, against $\lambda>2$.~\qed\end{proof}
   \step{s5512t3}{\case{$t\ge3$.}}
   \lean{Computations.C14.lambda_ge_nineteen}
   \begin{proof}\pf\ Then $2t(t-1)>0$, so $\lambda>2t^2\ge18$ by \stepref{s5512lam}, against
    $\lambda<19$.~\qed\end{proof}
   \step{s5512t2}{\case{$t=2$.}}
   \lean{Computations.C14.t2_parameters}
   \begin{proof}
    \step{s5512s}{$s\in\{3,5\}$.}
    \lean{Computations.C14.eq_three_or_five}
    \begin{proof}\pf\ By (iii), $(s-1)\mid4$ and $s\ge2t-1=3$ (computation C14).~\qed\end{proof}
    \step{s5512s3}{\case{$s=3$.}}
    \begin{proof}\pf\ Then $v=6=k$ by (iii), against $k<v$.~\qed\end{proof}
    \step{s5512s5}{\case{$s=5$.}}
    \begin{proof}\pf\ Then $(v,k,b,r,\lambda)=(15,10,21,14,9)$ by (iii), and
     Corollary~\hyperref[hc:cor:complement]{(the complementary parameters)} of the Hall--Connor
     section gives \textsc{False}.~\qed\end{proof}
    \qedstep
    \begin{proof}\pf\ The cases \stepref{s5512s3} and \stepref{s5512s5} are exhaustive by
     \stepref{s5512s}.~\qed\end{proof}
   \end{proof}
   \qedstep
   \begin{proof}\pf\ $t\ge1$ by (iii), so the cases \stepref{s5512t1}, \stepref{s5512t3} and
    \stepref{s5512t2} are exhaustive.~\qed\end{proof}
  \end{proof}
  \qedstep
  \begin{proof}\pf\ Theorem~\ref{thm:5.8} applies by \stepref{s5512fisher} and \asm{1}, so the
   cases are exhaustive.~\qed\end{proof}
 \end{proof}
 \qedstep
 \begin{proof}\pf\ Step \stepref{s5512mate} is Definition~\ref{def:flag}.~\qed\end{proof}
\end{proof}

\begin{remark}
 \lean{Computations.TrivialDesignException.design}
 \lean{Computations.TrivialDesignException.incidenceSignless}
 \lean{Computations.TrivialDesignException.completeAdjacency}
 \lean{Computations.TrivialDesignException.incidence_spectrum_checked}
 \lean{Computations.TrivialDesignException.complete_spectrum_checked}
 \lean{Computations.TrivialDesignException.incidenceIndex}
 \lean{Computations.TrivialDesignException.incidence_signless_eq}
 \lean{Computations.TrivialDesignException.spectrum_incidence_signless}
 \lean{Computations.TrivialDesignException.spectrum_complete}
 \lean{Computations.TrivialDesignException.spectrum_triangular}
 \lean{Computations.TrivialDesignException.spectrum_flagGraph}
 \lean{Computations.TrivialDesignException.spectrum_double_lineGraph}
 \lean{Computations.TrivialDesignException.flagGraph_cospectral_double_lineGraph}
 \lean{Computations.TrivialDesignException.lineGraph_completeBipartite_cospectral_double_lineGraph}
 \lean{Computations.TrivialDesignException.flagGraph_not_nonempty_iso_double_lineGraph}
 \lean{Computations.TrivialDesignException.not_flagGraphCharacterized}
 \lean{SimpleGraph.spectrumMultiset_lineGraph_of_card_le}
 The paper adds that for $t=2$ and $s=3$ the parameters are those of the trivial design with
 $(v,k,b,r)=(6,6,10,10)$, and that $L(L(K_6))$ provides an exception. The exception is genuine,
 and we prove it in \path{LineGraphs/Computations/TrivialDesignException.lean}. Let $D_1$ be the
 design $(6,6,10,10,10)$, whose ten blocks all equal its set of six points
 (\path{Computations.TrivialDesignException.design}); its incidence graph is $K_{10,6}$, so
 $F(D_1)\cong L(K_{10,6})$. For a graph $Y$ with $n$ vertices and $m\ge n$ edges, the spectrum of
 $L(Y)$ is the spectrum of the signless Laplacian of $Y$ shifted by $-2$, together with the
 eigenvalue $-2$ of multiplicity $m-n$
 (\path{SimpleGraph.spectrumMultiset_lineGraph_of_card_le}). Certificates for the signless
 Laplacian of $K_{10,6}$, of order 16, and for $A(K_6)$, of order 6, give through this identity
 the spectrum $14^1\,8^5\,4^9\,(-2)^{45}$ for both $F(D_1)$ and $L(L(K_6))=L(T(6))$
 (\path{Computations.TrivialDesignException.flagGraph_cospectral_double_lineGraph}). Thus
 $L(K_{10,6})$ and $L(T(6))$ are cospectral, the pair behind Errors 14, 16 and 17 (errata). They
 are not isomorphic, and $F(D_1)$ is not characterized by its spectrum in the sense of
 Definition~\ref{def:flag}
 (\path{Computations.TrivialDesignException.not_flagGraphCharacterized}): the flag graph of a
 design with replication number 10 has a clique of 10 vertices, the flags through one point,
 while a clique of $L(T(6))$ consists of edges of $T(6)$ through one vertex or of the three edges
 of a triangle, and so has at most 8 vertices, since $T(6)$ is 8-regular.
\end{remark}

\medskip\noindent\textbf{Example 5.13.}

\begin{cgss}{Example}{5.13}\label{ex:5.13}
 \lean{Design.not_flagGraphCharacterized_of_resolvable}
 \uses{def:flag,hc:def:design,hc:def:matrices,lem:design,lem:incidence,lem:rootgraph,lem:5.1,lem:5.3,thm:5.8,input:ex513,facts:sec5,st:2.2,st:2.6,st:3.1,st:3.2,st:3.5,st:6.1,st:6.2,st:6.6}
 \assume{\begin{pfenum}
  \item $k\ge2$ is an integer, and $D'$ is a design with parameters
   $(k(2k-1),k,4k^2-1,2k+1,1)$ on $P'$ and $\cB'$, with incidence matrix $N'$.
  \item $\pi:\cB'\to\{1,\dots,2k+1\}$ is a \emph{resolution} of $D'$: for each $i$, every point
   lies in exactly one block $B$ with $\pi(B)=i$.
  \item $H$ is the graph on $\cB'$ in which $B\ne B'$ are adjacent when $B$ and $B'$ are disjoint
   and $\pi(B)\ne\pi(B')$.
  \item $D_1$ is a design with parameters $(v_1,k_1,b_1,r_1,\lambda_1)=
   (k(2k-1),(k-1)(2k-1),k(2k+1),(k-1)(2k+1),k(2k-3))$.
 \end{pfenum}}
 \prove{\begin{pfenum}
  \item $H$ is connected and $\kappa$-regular on $4k^2-1$ vertices, where $\kappa=2k(k-1)$, and
   $\Spec(H)=\ms{\kappa^{1},(-(k-1))^{2k},k^{\ell},(-k)^{\ell}}$, where $\ell=(k-1)(2k+1)$.
   This is the spectrum of case (iii) of Theorem~\ref{thm:5.8} with $s=2k-1$ and $t=k-1$.
  \item $L(H)$ is cospectral with $F(D_1)$.
  \item $F(D_1)$ is not characterized by its spectrum.
 \end{pfenum}}
\end{cgss}
\begin{proof}
 \pfsketch\ \repair\ The paper computes the spectrum of $H$ and leaves the rest to the reader;
 ``not characterized'' is meant in the sense of Definition~\ref{def:flag} (errata, Implicit
 hypothesis 6), and then it is EX4. We will split $\re^{\cB'}$ into four eigenspaces of $A(H)$,
 compare the signless Laplacians of $H$ and of $I(D_1)$ as in Lemma 5.5, and then apply EX4. With
 $s=2k-1$, $t=k-1$ we have $(s+1)t=\kappa$, $s(s+2)=4k^2-1$, $c=k^2$ and
 $\frac12(s-1)(s+2)=\ell$ in Theorem 5.8 (iii).
 \pflet{$c_i\in\re^{\cB'}$ be the indicator vector of $\pi^{-1}(i)$ for $1\le i\le2k+1$, $R$ the
  $\cB'\times\cB'$ matrix with $R(B,B')=1$ when $B\ne B'$ and $\pi(B)=\pi(B')$ and $0$ otherwise,
  and $J$ the all-ones $\cB'\times\cB'$ matrix.}
 \step{s5exclass}{$N'c_i=\Zj$, each $\pi^{-1}(i)$ has $2k-1$ blocks, $\sum_ic_i=\Zj$ and
  $R=\sum_ic_ic_i^T-I$.}
 \begin{proof}\pf\ By \asm[2]{0}, the blocks of $\pi^{-1}(i)$ partition $P'$; each has $k$
  points and $|P'|=k(2k-1)$. Every block lies in exactly one class.~\qed\end{proof}
 \step{s5exadj}{$A(H)=J-I-R-(N'^TN'-kI)$.}
 \begin{proof}
  \step{s5exadjd}{$(N'^TN')(B,B')=|D'(B)\cap D'(B')|$, and $(N'^TN')(B,B)=k$.}
  \begin{proof}\pf\ The entry is the number of points on both blocks, and every block has $k$
   points by \asm[1]{0}.~\qed\end{proof}
  \step{s5exadj01}{$(N'^TN')(B,B')\in\{0,1\}$ for $B\ne B'$.}
  \begin{proof}\pf\ Two distinct points lie in exactly $\lambda=1$ common block, by \asm[1]{0},
   so distinct blocks share at most one point; apply \stepref{s5exadjd}.~\qed\end{proof}
  \step{s5exadjc}{$(N'^TN')(B,B')=0$ for $B\ne B'$ with $\pi(B)=\pi(B')$.}
  \begin{proof}\pf\ The blocks of $\pi^{-1}(\pi(B))$ are disjoint by \stepref{s5exclass}; apply
   \stepref{s5exadjd}.~\qed\end{proof}
  \qedstep
  \begin{proof}\pf\ By Steps \stepref{s5exadjd} to \stepref{s5exadjc}, the matrix
   $J-I-R-(N'^TN'-kI)$ has $(B,B')$ entry $0$ when $B=B'$ or $\pi(B)=\pi(B')$, and
   $1-|D'(B)\cap D'(B')|$ otherwise, which is $1$ exactly when $B$ and $B'$ are disjoint. Compare
   the entries with \asm[3]{0}.~\qed\end{proof}
 \end{proof}
 \step{s5exgram}{$N'N'^T=2kI+J$, and $\rk N'=k(2k-1)$.}
 \begin{proof}\pf\ Lemma~\ref{lem:design}, assertion 3, with $r-\lambda=2k$; its spectrum is
  positive, so $\rk(N'N'^T)=k(2k-1)$ by \ST{Theorem 3.2}, and $\rk N'=\rk(N'N'^T)$ by
  \ST{Theorem 3.5}.~\qed\end{proof}
 \pflet{$W_1=\{\sum_ia_ic_i:\sum_ia_i=0\}$, $W_2=\{N'^T\Zy:\Zy^T\Zj=0\}$ and
  $W_3=\{\Zx:N'\Zx=0,\ c_i^T\Zx=0\text{ for all }i\}$.}
 \step{s5exvalues}{$A(H)\Zj=\kappa\Zj$, and $A(H)$ acts on $W_1$, $W_2$, $W_3$ as $-(k-1)$, $-k$
  and $k$.}
 \begin{proof}
  \step{s5exvj}{$A(H)\Zj=\kappa\Zj$.}
  \begin{proof}\pf\ $J\Zj=(4k^2-1)\Zj$, $R\Zj=(2k-2)\Zj$ by \stepref{s5exclass}, and
   $N'^TN'\Zj=(2k+1)k\Zj$ from the row and column sums of $N'$; by \stepref{s5exadj},
   $A(H)\Zj=(4k^2-2-(2k-2)-(2k+1)k+k)\Zj=\kappa\Zj$.~\qed\end{proof}
  \step{s5exv1}{$A(H)\Zx=-(k-1)\Zx$ for $\Zx\in W_1$.}
  \begin{proof}\pf\ By \stepref{s5exclass}, $J\Zx=(2k-1)(\sum_ia_i)\Zj=0$,
   $R\Zx=(2k-2)\Zx$ and $N'\Zx=(\sum_ia_i)\Zj=0$; apply \stepref{s5exadj}.~\qed\end{proof}
  \step{s5exv2}{$A(H)\Zx=-k\Zx$ for $\Zx\in W_2$.}
  \begin{proof}\pf\ For $\Zx=N'^T\Zy$ with $\Zy^T\Zj=0$, $c_i^T\Zx=\Zj^T\Zy=0$ by
   \stepref{s5exclass}, so $J\Zx=0$ and $R\Zx=-\Zx$; and $N'^TN'\Zx=N'^T(2kI+J)\Zy=2k\Zx$ by
   \stepref{s5exgram}. Apply \stepref{s5exadj}.~\qed\end{proof}
  \step{s5exv3}{$A(H)\Zx=k\Zx$ for $\Zx\in W_3$.}
  \begin{proof}\pf\ $J\Zx=0$ and $R\Zx=-\Zx$ by \stepref{s5exclass}, and $N'^TN'\Zx=0$; apply
   \stepref{s5exadj}.~\qed\end{proof}
  \qedstep
  \begin{proof}\pf\ Steps \stepref{s5exvj} to \stepref{s5exv3}.~\qed\end{proof}
 \end{proof}
 \step{s5exdims}{$\dim W_1=2k$ and $\dim W_2=\dim W_3=\ell$.}
 \begin{proof}
  \step{s5exd12}{$\dim W_1=2k$ and $\dim W_2=\ell$.}
  \begin{proof}\pf\ By \stepref{s5exclass}, the $2k+1$ vectors $c_i$ have disjoint nonempty
   supports, so they are independent, and $W_1$ is a hyperplane of their span. $N'^T$ is
   injective by \stepref{s5exgram}, so $\dim W_2=k(2k-1)-1=\ell$.~\qed\end{proof}
  \step{s5exd3}{$\dim W_3=\ell$.}
  \begin{proof}
   \step{s5exd3k}{$\dim\ker N'=(4k^2-1)-k(2k-1)$.}
   \begin{proof}\pf\ $N'$ has $4k^2-1$ columns and rank $k(2k-1)$ by
    \stepref{s5exgram}.~\qed\end{proof}
   \step{s5exd3w}{$W_1\sbs\ker N'$.}
   \begin{proof}\pf\ $N'\sum_ia_ic_i=(\sum_ia_i)\Zj$ by \stepref{s5exclass}.~\qed\end{proof}
   \step{s5exd3p}{$W_3$ is the orthogonal complement of $W_1$ in $\ker N'$.}
   \begin{proof}\pf\ A vector $\Zx$ is orthogonal to $W_1$ exactly when all $c_i^T\Zx$ are
    equal, so $W_3$ lies in the complement. Conversely, if $\Zx\in\ker N'$ and $\Zx\perp W_1$,
    then all $c_i^T\Zx$ are equal, and their sum is $\Zj^T\Zx=\frac1k\Zj^TN'\Zx=0$ by
    \stepref{s5exclass}, since the column sums of $N'$ are $k$; so $\Zx\in
    W_3$.~\qed\end{proof}
   \qedstep
   \begin{proof}\pf\ By Steps \stepref{s5exd3k} to \stepref{s5exd3p} and \stepref{s5exd12},
    $\dim W_3=(4k^2-1)-k(2k-1)-2k=\ell$.~\qed\end{proof}
  \end{proof}
  \qedstep
  \begin{proof}\pf\ Steps \stepref{s5exd12} and \stepref{s5exd3}.~\qed\end{proof}
 \end{proof}
 \step{s5exspec}{Assertion 1 holds.}
 \begin{proof}
  \step{s5exspecH}{$\Spec(H)=\ms{\kappa^{1},(-(k-1))^{2k},k^{\ell},(-k)^{\ell}}$.}
  \begin{proof}\pf\ The values $\kappa$, $-(k-1)$, $-k$, $k$ are distinct since $k\ge2$, so by
   \ST{Theorem 2.2}, \stepref{s5exvalues} and \stepref{s5exdims}, their multiplicities are at
   least $1$, $2k$, $\ell$, $\ell$. These add up to $4k^2-1$, the order of $A(H)$, so by
   \ST{Theorem 2.2} they are the multiplicities and there is no other eigenvalue.~\qed\end{proof}
  \step{s5exspec58}{$\Spec(H)$ is the spectrum of case (iii) of Theorem~\ref{thm:5.8} with
   $s=2k-1$ and $t=k-1$.}
  \begin{proof}\pf\ With these values, $(s+1)t=2k(k-1)=\kappa$, $s+1=2k$, $-t=-(k-1)$,
   $c=t(s+1)(s-t)/(s-1)=(k-1)\cdot2k\cdot k/(2k-2)=k^2$, so $\sqrt c=k$, and
   $\frac12(s-1)(s+2)=(k-1)(2k+1)=\ell$. Compare with \stepref{s5exspecH}.~\qed\end{proof}
  \qedstep
  \begin{proof}\pf\ $H$ is $\kappa$-regular on $|\cB'|=4k^2-1$ vertices by \stepref{s5exvalues}, and
   connected by Lemma~\ref{lem:5.1}, since $\mult_H(\kappa)=1$ by
   \stepref{s5exspecH}. The spectrum is given by Steps \stepref{s5exspecH} and
   \stepref{s5exspec58}.~\qed\end{proof}
 \end{proof}
 \step{s5exline}{Assertion 2 holds.}
 \begin{proof}
  \pflet{$M_H$ and $M_1$ be the incidence matrices of $H$ and of $I(D_1)$.}
  \step{s5exCH}{$\Spec(M_HM_H^T)=\ms{(2\kappa)^{1},(\kappa-k+1)^{2k},(\kappa+k)^{\ell},
   (\kappa-k)^{\ell}}$, with all members positive.}
  \begin{proof}\pf\ $M_HM_H^T=\kappa I+A(H)$ by \ST{Theorem 6.2} and \stepref{s5exspec}; shift
   $\Spec(H)$ by $\kappa$ with \ST{Theorem 3.1}. $\kappa-k=k(2k-3)>0$ as $k\ge2$.~\qed\end{proof}
  \step{s5exC1}{$\Spec(M_1M_1^T)$ is $\Spec(M_HM_H^T)$ with one copy of $0$ added.}
  \begin{proof}
   \step{s5exC1a}{$\Spec(N_1N_1^T)=\ms{(r_1k_1)^{1},(2k-1)^{v_1-1}}$ and $\rk N_1=v_1$, for the
    incidence matrix $N_1$ of $D_1$.}
   \begin{proof}\pf\ Lemma~\ref{lem:design}, assertion 3, with $r_1-\lambda_1=2k-1>0$; then
    \ST{Theorem 3.2} and \ST{Theorem 3.5}.~\qed\end{proof}
   \step{s5exC1b}{$\Spec(M_1M_1^T)$ consists of $0$ and $r_1+k_1$, of
    $\frac12\bigl(r_1+k_1\pm\sqrt{(r_1-k_1)^2+4(2k-1)}\bigr)$ with multiplicity $v_1-1$ each, and
    of $k_1$ with multiplicity $b_1-v_1$.}
   \begin{proof}\pf\ Lemma~\ref{lem:rootgraph}, assertion 3, for $Z=I(D_1)$, which applies by
    Lemma~\ref{lem:incidence}; then \ST{Theorem 6.6} with $s=r_1$, $t=k_1$, $\rho=v_1$ and
    \stepref{s5exC1a}. The roots for $r_1k_1$ are $0$ and $r_1+k_1$.~\qed\end{proof}
   \qedstep
   \begin{proof}\pf\ Substitute \asm[4]{0}: $r_1+k_1=2\kappa$, $r_1-k_1=2(k-1)$, so the square
    root is $2k$ and the pair is $\kappa\pm k$; $v_1-1=\ell$, $k_1=\kappa-k+1$ and
    $b_1-v_1=2k$. Compare with \stepref{s5exCH}.~\qed\end{proof}
  \end{proof}
  \step{s5exedges}{$|E(H)|=|E(I(D_1))|=v_1r_1$.}
  \begin{proof}\pf\ $|E(H)|=\frac12\kappa(4k^2-1)=k(k-1)(2k-1)(2k+1)$ by
   \path{SimpleGraph.sum_degrees_eq_twice_card_edges} and \stepref{s5exspec}, and this is
   $v_1r_1$, which is $|E(I(D_1))|$ by Lemma~\ref{lem:incidence}.~\qed\end{proof}
  \step{s5exnz}{$\mult_{2I+A(L(H))}(\gamma)=\mult_{2I+A(F(D_1))}(\gamma)$ for every
   $\gamma\ne0$.}
  \begin{proof}\pf\ By \ST{Theorem 6.2}, $M_H^TM_H=2I+A(L(H))$ and $M_1^TM_1=2I+A(F(D_1))$; by
   \ST{Theorem 6.1}, their multiplicities at $\gamma\ne0$ are those of $M_HM_H^T$ and
   $M_1M_1^T$, which agree by \stepref{s5exC1}.~\qed\end{proof}
  \qedstep
  \begin{proof}\pf\ Both matrices have order $v_1r_1$ by \stepref{s5exedges}, so their
   multiplicities at $0$ agree by \ST{Theorem 2.2} and \stepref{s5exnz}. Thus
   $\Spec(2I+A(L(H)))=\Spec(2I+A(F(D_1)))$ by \ST{Proposition 2.6}; shift by $-2$ with
   \ST{Theorem 3.1}, and apply \ST{Proposition 2.6}.~\qed\end{proof}
 \end{proof}
 \step{s5exnot}{Assertion 3 holds.}
 \begin{proof}
  \step{s5exnonbip}{$H$ is not bipartite.}
  \begin{proof}\pf\ $H$ is connected and $\kappa$-regular, and $-\kappa\notin\Spec(H)$ by
   \stepref{s5exspec}, since $\kappa>k$ for $k\ge2$; apply Lemma~\ref{lem:5.3}.~\qed\end{proof}
  \step{s5exnoiso}{\assume{$D_2$ is a design with the parameters of $D_1$, and
   $L(H)\cong F(D_2)$.}\prove{\textsc{False}.}}
  \begin{proof}\pf\ $F(D_2)=L(I(D_2))$, and $I(D_2)$ is bipartite without isolated vertices by
   Lemma~\ref{lem:incidence}. $H$ is connected on $4k^2-1\ge15$ vertices by \stepref{s5exspec},
   and not bipartite by \stepref{s5exnonbip}. This contradicts EX4 of the
   \hyperref[input:ex513]{inputs for Example 5.13}. \new\ The Lean proof,
   \path{Design.not_nonempty_iso_flagGraph_lineGraph_of_degree}, uses a clique bound in place of
   EX4, so the example does not need EX4: the $r_1$ flags on a point of $D_2$ form an
   $r_1$-clique of $F(D_2)$ (\path{Design.exists_clique_flagGraph_card_replication}), while
   every clique of $L(H)$ has at most $\kappa$ vertices, since $H$ is $\kappa$-regular with
   $\kappa\ge3$ (\path{SimpleGraph.card_clique_lineGraph_le}); and
   $r_1=(k-1)(2k+1)=\kappa+k-1>\kappa=2k(k-1)\ge4$ for $k\ge2$.~\qed\end{proof}
  \qedstep
  \begin{proof}\pf\ $L(H)$ is cospectral with $F(D_1)$ by \stepref{s5exline}, and is isomorphic
   to no $F(D_2)$ by \stepref{s5exnoiso}; this is the negation of Definition~\ref{def:flag}
   for $D_1$.~\qed\end{proof}
 \end{proof}
 \qedstep
 \begin{proof}\pf\ Steps \stepref{s5exspec}, \stepref{s5exline} and \stepref{s5exnot}.~\qed\end{proof}
\end{proof}

\begin{remark}
 \lean{Computations.C15.flagGraph_cospectral_double_lineGraph}
 \lean{Computations.C15.flagGraph_not_nonempty_iso_double_lineGraph}
 \lean{Computations.C15.design}\lean{Computations.C15.petersen}
 \lean{Computations.C15.petersenIso}\lean{Computations.C15.nonempty_iso_petersen}
 \lean{Computations.C15.four_cliques_graph1}\lean{Computations.C15.four_cliques_graph2}
 By EX1 and EX2, the resolvable designs of the example exist for $k=2^d$ with $d\ge1$ and for
 $k=3$, and by EX3 the designs $D_1$ exist for $k=2$ and $k=4$, with parameters $(6,3,10,5,2)$
 and $(28,21,36,27,20)$; for these two, the flag graph is not characterized by its spectrum. For
 $k=2$ the paper describes $H$ as the line graph of the Petersen graph; we do not check this.
 Computation C15 checks that the flag graph of the design $(6,3,10,5,2)$ is cospectral with
 $L(L(\text{Petersen}))$ and not isomorphic to it; the two graphs have 30 and 15 cliques of
 size 4. In Lean the flag graph is the line graph of the incidence graph of
 \path{LineGraphs/Computations/C13Design.lean}, and the Petersen graph is the settled
 \path{SimpleGraph.petersen}. The certificate is computed on \path{Computations.C15.petersen},
 the pentagon, the pentagram and five spokes, and \path{Computations.C15.petersenIso} is an
 explicit isomorphism from it to \path{SimpleGraph.petersen}, through which the two statements
 are transported. For $k=3$ the parameters of $D_1$ are $(15,10,21,14,9)$, and no such design exists, by the Hall--Connor section, as in the proof of
 Remark 5.12. Computation C23 (\path{LineGraphs/Computations/C23.lean}, with the data of
 \path{cgss/certificates/c23_lean_data.py}) builds the designs of the instances $k=2$ and $k=4$ and proves
 both instances, and the case $k=3$ above, in Lean; the existence of the designs of EX1 for every
 $k=2^d$ is a citation, outside the scope.
\end{remark}

\endgroup
